% Recovered from the previously completed PDF; no new mathematical content.
% Prose is editable; formula boxes retain original symbol positions.
% Build from repository root: xelatex -no-shell-escape -halt-on-error -output-directory=... print/fudan/NAME.tex
\documentclass[11pt,a4paper]{article}
\usepackage[margin=23mm,headheight=14pt,headsep=8.5mm,footskip=10mm]{geometry}
\usepackage{amsmath,amssymb}
\DeclareMathAlphabet{\mathsf}{OT1}{cmss}{m}{n}
\usepackage[no-math]{fontspec}
\usepackage{xeCJK}
\setmainfont{Latin Modern Roman}
\setCJKmainfont{FandolSong-Regular.otf}[Path=assets/fonts/fandol/,BoldFont=FandolSong-Bold.otf,ItalicFont=FandolSong-Regular.otf,Script=CustomDefault]
\xeCJKsetup{PunctStyle=quanjiao,CheckSingle=true}
\usepackage{setspace}\setstretch{1.12}
\setlength{\parindent}{2em}\setlength{\parskip}{3pt}
\emergencystretch=6em
\hfuzz=9pt
\usepackage{titlesec}
\titleformat{\section}{\centering\normalfont\Large\bfseries}{}{0pt}{}
\titlespacing*{\section}{0pt}{1.8\baselineskip}{1.0\baselineskip}
\usepackage{fancyhdr}\pagestyle{fancy}\fancyhf{}
\fancyhead[L]{\small 数学手记·高等代数}\fancyhead[R]{\small 14A07–14A13}
\fancyfoot[C]{\thepage}\renewcommand{\headrulewidth}{0.4pt}
\usepackage[hidelinks,unicode]{hyperref}
\hypersetup{pdftitle={14A07–14A13：题目、解答与小结},pdfauthor={Vincent Arqain}}
\newcommand{\RB}[4]{\raisebox{0pt}[#2bp][#3bp]{\hbox to #1bp{#4\hss}}}
\newcommand{\RG}[5]{\hbox to 0pt{\kern#1bp\raisebox{#2bp}[0pt][0pt]{\fontsize{#3bp}{#3bp}\csname font#4\endcsname\char#5}\hss}}
\newcommand{\RT}[4]{\hbox to 0pt{\kern#1bp\raisebox{#2bp}[0pt][0pt]{\fontsize{#3bp}{#3bp}\selectfont#4}\hss}}
\newcommand{\RS}[4]{\hbox to 0pt{\kern#1bp\raisebox{#2bp}[0pt][0pt]{\fontsize{#3bp}{#3bp}\selectfont\ensuremath{#4}}\hss}}
\newcommand{\RR}[4]{\hbox to 0pt{\kern#1bp\raisebox{#2bp}[0pt][0pt]{\rule{#3bp}{#4bp}}\hss}}
\newcommand{\fontmi}{\usefont{OML}{cmm}{m}{it}}
\newcommand{\fontsy}{\usefont{OMS}{cmsy}{m}{n}}
\newcommand{\fontex}{\usefont{OMX}{cmex}{m}{n}}
\newcommand{\fontrm}{\usefont{OT1}{cmr}{m}{n}}
\newcommand{\fontbf}{\usefont{OT1}{cmr}{bx}{n}}
\newcommand{\fontam}{\usefont{U}{msa}{m}{n}}
\newcommand{\fontbb}{\usefont{U}{msb}{m}{n}}
\widowpenalty=10000\clubpenalty=10000\raggedbottom
\begin{document}
\begin{center}{\LARGE\bfseries 14A07–14A13\par}\vspace{8pt}复旦大学高等代数习题集·题目与解答\end{center}
% Original page 1.
题目、简洁解答与方法总结

\[
\RB{331.118}{7.538}{23.142}{%
\RT{0.000}{0.004}{7.9701}{题}\RT{8.727}{0.004}{7.9701}{目}\RT{17.454}{0.004}{7.9701}{来}\RT{26.181}{0.004}{7.9701}{源}\RT{34.908}{0.004}{7.9701}{：}\RT{43.635}{0.004}{7.9701}{谢}\RT{52.362}{0.004}{7.9701}{启}\RT{61.089}{0.004}{7.9701}{鸿}\RT{69.816}{0.004}{7.9701}{《}\RT{78.543}{0.004}{7.9701}{复}\RT{87.269}{0.004}{7.9701}{旦}\RT{95.996}{0.004}{7.9701}{大}\RT{104.723}{0.004}{7.9701}{学}\RT{113.450}{0.004}{7.9701}{高}\RT{122.177}{0.004}{7.9701}{等}\RT{130.904}{0.004}{7.9701}{代}\RT{139.631}{0.004}{7.9701}{数}\RT{148.358}{0.004}{7.9701}{习}\RT{157.085}{0.004}{7.9701}{题}\RT{165.812}{0.004}{7.9701}{集}\RT{174.539}{0.004}{7.9701}{》}\RT{178.902}{0.004}{7.9701}{（}\RT{187.629}{0.004}{7.9701}{2}\RT{192.263}{0.004}{7.9701}{0}\RT{196.897}{0.004}{7.9701}{2}\RT{201.531}{0.004}{7.9701}{3}\RT{209.255}{0.004}{7.9701}{版}\RT{217.982}{0.004}{7.9701}{）}\RT{222.354}{0.004}{7.9701}{，}\RT{231.073}{0.004}{7.9701}{§}\RT{237.059}{0.004}{7.9701}{1}\RT{241.693}{0.004}{7.9701}{.}\RT{244.268}{0.004}{7.9701}{2}\RT{248.907}{0.004}{7.9701}{，}\RT{257.634}{0.004}{7.9701}{原}\RT{266.361}{0.004}{7.9701}{书}\RT{275.088}{0.004}{7.9701}{第}\RT{286.905}{0.004}{7.9701}{1}\RT{291.539}{0.004}{7.9701}{2}\RT{296.173}{0.004}{7.9701}{–}\RT{300.807}{0.004}{7.9701}{1}\RT{305.441}{0.004}{7.9701}{3}\RT{313.164}{0.004}{7.9701}{页}\RT{321.891}{0.004}{7.9701}{。}\RG{33.126}{-17.816}{7.9701}{sy}{48}\RT{55.374}{-17.816}{7.9701}{T}\RT{0.000}{-20.896}{7.9701}{记}\RT{8.727}{-20.896}{7.9701}{号}\RT{17.454}{-20.896}{7.9701}{：}\RG{26.181}{-20.896}{7.9701}{mi}{65}\RG{38.641}{-20.896}{7.9701}{rm}{61}\RG{48.429}{-20.896}{7.9701}{mi}{65}\RT{63.761}{-20.896}{7.9701}{表}\RT{72.488}{-20.896}{7.9701}{示}\RT{81.215}{-20.896}{7.9701}{转}\RT{89.941}{-20.896}{7.9701}{置}\RT{98.668}{-20.896}{7.9701}{；}\RT{107.396}{-20.896}{7.9701}{I}\RT{110.738}{-20.896}{7.9701}{m}\RG{120.007}{-20.896}{7.9701}{mi}{65}\RT{130.041}{-20.896}{7.9701}{与}\RT{141.858}{-20.896}{7.9701}{k}\RT{146.753}{-20.896}{7.9701}{e}\RT{150.873}{-20.896}{7.9701}{r}\RG{156.031}{-20.896}{7.9701}{mi}{65}\RT{166.066}{-20.896}{7.9701}{分}\RT{174.793}{-20.896}{7.9701}{别}\RT{183.520}{-20.896}{7.9701}{表}\RT{192.247}{-20.896}{7.9701}{示}\RT{200.974}{-20.896}{7.9701}{像}\RT{209.701}{-20.896}{7.9701}{空}\RT{218.428}{-20.896}{7.9701}{间}\RT{227.155}{-20.896}{7.9701}{与}\RT{235.882}{-20.896}{7.9701}{核}\RT{244.608}{-20.896}{7.9701}{空}\RT{253.335}{-20.896}{7.9701}{间}\RT{262.062}{-20.896}{7.9701}{。}%
}
\]
\section*{2014A07递推向量与不可约多项式}
题目设\allowbreak{}\hskip.1em \RB{8.365}{8.312}{0.500}{%
\RG{0.000}{0.001}{10.9091}{mi}{65}%
}\hskip.1em\allowbreak{} 是有理数域\allowbreak{}\hskip.1em \RB{8.627}{8.181}{2.452}{%
\RG{0.000}{0.001}{10.9091}{bb}{81}%
}\hskip.1em\allowbreak{} 上的\allowbreak{}\hskip.1em \RB{5.638}{7.886}{0.500}{%
\RG{0.000}{0.001}{10.9091}{rm}{52}%
}\hskip.1em\allowbreak{} 阶方阵，\allowbreak{}\hskip.1em \RB{61.679}{5.323}{2.604}{%
\RG{0.000}{0.001}{10.9091}{mi}{11}\RG{6.978}{-1.635}{7.9701}{rm}{49}\RG{11.872}{0.001}{10.9091}{mi}{59}\RG{16.704}{0.001}{10.9091}{mi}{11}\RG{23.698}{-1.635}{7.9701}{rm}{50}\RG{28.592}{0.001}{10.9091}{mi}{59}\RG{33.424}{0.001}{10.9091}{mi}{11}\RG{40.418}{-1.635}{7.9701}{rm}{51}\RG{45.313}{0.001}{10.9091}{mi}{59}\RG{50.145}{0.001}{10.9091}{mi}{11}\RG{57.139}{-1.635}{7.9701}{rm}{52}%
}\hskip.1em\allowbreak{} 是\allowbreak{}\hskip.1em \RB{8.627}{8.181}{2.452}{%
\RG{0.000}{0.001}{10.9091}{bb}{81}%
}\hskip.1em\allowbreak{} 上的\allowbreak{}\hskip.1em \RB{5.638}{7.886}{0.500}{%
\RG{0.000}{0.001}{10.9091}{rm}{52}%
}\hskip.1em\allowbreak{} 维列向量，满足

\[
\RB{330.755}{8.312}{2.604}{%
\RG{0.000}{0.001}{10.9091}{mi}{65}\RG{8.181}{0.001}{10.9091}{mi}{11}\RG{23.084}{0.001}{10.9091}{rm}{61}\RG{34.599}{0.001}{10.9091}{mi}{11}\RG{46.470}{0.001}{10.9091}{mi}{59}\RG{62.211}{0.001}{10.9091}{mi}{65}\RG{70.393}{0.001}{10.9091}{mi}{11}\RG{85.312}{0.001}{10.9091}{rm}{61}\RG{96.826}{0.001}{10.9091}{mi}{11}\RG{108.698}{0.001}{10.9091}{mi}{59}\RG{124.439}{0.001}{10.9091}{mi}{65}\RG{132.620}{0.001}{10.9091}{mi}{11}\RG{147.538}{0.001}{10.9091}{rm}{61}\RG{159.053}{0.001}{10.9091}{mi}{11}\RG{170.925}{0.001}{10.9091}{mi}{59}\RG{197.575}{0.001}{10.9091}{mi}{65}\RG{205.757}{0.001}{10.9091}{mi}{11}\RG{220.675}{0.001}{10.9091}{rm}{61}\RG{232.190}{0.001}{10.9091}{sy}{0}\RG{240.674}{0.001}{10.9091}{mi}{11}\RG{254.970}{0.001}{10.9091}{sy}{0}\RG{265.879}{0.001}{10.9091}{mi}{11}\RG{280.175}{0.001}{10.9091}{sy}{0}\RG{291.084}{0.001}{10.9091}{mi}{11}\RG{305.380}{0.001}{10.9091}{sy}{0}\RG{316.289}{0.001}{10.9091}{mi}{11}\RG{328.161}{0.001}{10.9091}{mi}{58}\RG{15.160}{-1.635}{7.9701}{rm}{49}\RG{41.577}{-1.635}{7.9701}{rm}{50}\RG{77.387}{-1.635}{7.9701}{rm}{50}\RG{103.804}{-1.635}{7.9701}{rm}{51}\RG{139.615}{-1.635}{7.9701}{rm}{51}\RG{166.032}{-1.635}{7.9701}{rm}{52}\RG{212.751}{-1.635}{7.9701}{rm}{52}\RG{247.653}{-1.635}{7.9701}{rm}{49}\RG{272.858}{-1.635}{7.9701}{rm}{50}\RG{298.062}{-1.635}{7.9701}{rm}{51}\RG{323.267}{-1.635}{7.9701}{rm}{52}%
}
\]
证明：若\allowbreak{}\hskip.1em \RB{31.918}{8.313}{2.854}{%
\RG{0.000}{0.002}{10.9091}{mi}{11}\RG{6.978}{-1.635}{7.9701}{rm}{49}\RG{14.902}{0.002}{10.9091}{sy}{54}\RG{14.902}{0.002}{10.9091}{rm}{61}\RG{26.400}{0.002}{10.9091}{rm}{48}%
}\hskip.1em\allowbreak{} ，则\allowbreak{}\hskip.1em \RB{72.646}{8.684}{3.225}{%
\RG{0.000}{0.002}{10.9091}{sy}{102}\RG{5.454}{0.002}{10.9091}{mi}{11}\RG{12.432}{-1.635}{7.9701}{rm}{49}\RG{17.327}{0.002}{10.9091}{mi}{59}\RG{22.159}{0.002}{10.9091}{mi}{11}\RG{29.153}{-1.635}{7.9701}{rm}{50}\RG{34.047}{0.002}{10.9091}{mi}{59}\RG{38.879}{0.002}{10.9091}{mi}{11}\RG{45.873}{-1.635}{7.9701}{rm}{51}\RG{50.767}{0.002}{10.9091}{mi}{59}\RG{55.599}{0.002}{10.9091}{mi}{11}\RG{62.593}{-1.635}{7.9701}{rm}{52}\RG{67.488}{0.002}{10.9091}{sy}{103}%
}\hskip.1em\allowbreak{} 是\allowbreak{}\hskip.1em \RB{13.024}{9.607}{2.450}{%
\RG{0.000}{0.002}{10.9091}{bb}{81}\RG{8.485}{3.960}{7.9701}{rm}{52}%
}\hskip.1em\allowbreak{} 的一组基。

解：令

\[
\RB{124.275}{10.155}{3.223}{%
\RG{36.396}{4.508}{7.9701}{rm}{52}\RG{58.562}{4.508}{7.9701}{rm}{51}\RG{80.728}{4.508}{7.9701}{rm}{50}\RG{0.000}{0.004}{10.9091}{mi}{112}\RG{5.489}{0.004}{10.9091}{rm}{40}\RG{9.731}{0.004}{10.9091}{mi}{116}\RG{13.671}{0.004}{10.9091}{rm}{41}\RG{20.925}{0.004}{10.9091}{rm}{61}\RG{32.457}{0.004}{10.9091}{mi}{116}\RG{43.714}{0.004}{10.9091}{rm}{43}\RG{54.623}{0.004}{10.9091}{mi}{116}\RG{65.880}{0.004}{10.9091}{rm}{43}\RG{76.788}{0.004}{10.9091}{mi}{116}\RG{88.046}{0.004}{10.9091}{rm}{43}\RG{98.955}{0.004}{10.9091}{mi}{116}\RG{105.318}{0.004}{10.9091}{rm}{43}\RG{116.215}{0.004}{10.9091}{rm}{49}\RG{121.681}{0.004}{10.9091}{mi}{58}%
}
\]
由递推关系，\allowbreak{}\hskip.1em \RB{61.154}{9.528}{3.695}{%
\RG{0.000}{-0.001}{10.9091}{mi}{11}\RG{6.978}{-1.637}{7.9701}{mi}{106}\RG{14.605}{-0.001}{10.9091}{rm}{61}\RG{26.119}{-0.001}{10.9091}{mi}{65}\RG{34.301}{3.957}{7.9701}{mi}{106}\RG{38.352}{3.957}{7.9701}{sy}{0}\RG{45.169}{3.957}{7.9701}{rm}{49}\RG{50.064}{-0.001}{10.9091}{mi}{11}\RG{57.042}{-1.637}{7.9701}{rm}{49}%
}\hskip.1em\allowbreak{} （\allowbreak{}\hskip.1em \RB{45.299}{7.884}{2.737}{%
\RG{0.000}{-0.001}{10.9091}{rm}{49}\RG{8.485}{-0.001}{10.9091}{sy}{20}\RG{19.999}{-0.001}{10.9091}{mi}{106}\RG{28.146}{-0.001}{10.9091}{sy}{20}\RG{39.661}{-0.001}{10.9091}{rm}{52}%
}\hskip.1em\allowbreak{} ），且

\[
\RB{192.038}{8.680}{3.228}{%
\RG{0.000}{-0.001}{10.9091}{mi}{112}\RG{5.488}{-0.001}{10.9091}{rm}{40}\RG{9.731}{-0.001}{10.9091}{mi}{65}\RG{17.912}{-0.001}{10.9091}{rm}{41}\RG{22.154}{-0.001}{10.9091}{mi}{11}\RG{37.057}{-0.001}{10.9091}{rm}{61}\RG{48.571}{-0.001}{10.9091}{mi}{65}\RG{56.753}{-0.001}{10.9091}{mi}{11}\RG{71.049}{-0.001}{10.9091}{rm}{43}\RG{81.958}{-0.001}{10.9091}{mi}{11}\RG{96.254}{-0.001}{10.9091}{rm}{43}\RG{107.163}{-0.001}{10.9091}{mi}{11}\RG{121.458}{-0.001}{10.9091}{rm}{43}\RG{132.368}{-0.001}{10.9091}{mi}{11}\RG{146.663}{-0.001}{10.9091}{rm}{43}\RG{157.572}{-0.001}{10.9091}{mi}{11}\RG{172.475}{-0.001}{10.9091}{rm}{61}\RG{183.973}{-0.001}{10.9091}{rm}{48}\RG{189.443}{-0.001}{10.9091}{mi}{58}\RG{29.133}{-1.637}{7.9701}{rm}{49}\RG{63.731}{-1.637}{7.9701}{rm}{52}\RG{88.936}{-1.637}{7.9701}{rm}{52}\RG{114.141}{-1.637}{7.9701}{rm}{51}\RG{139.346}{-1.637}{7.9701}{rm}{50}\RG{164.551}{-1.637}{7.9701}{rm}{49}%
}
\]
先证明\allowbreak{}\hskip.1em \RB{17.323}{8.682}{3.227}{%
\RG{0.000}{0.000}{10.9091}{mi}{112}\RG{5.489}{0.000}{10.9091}{rm}{40}\RG{9.731}{0.000}{10.9091}{mi}{116}\RG{13.671}{0.000}{10.9091}{rm}{41}%
}\hskip.1em\allowbreak{} 在\allowbreak{}\hskip.1em \RB{8.627}{8.180}{2.453}{%
\RG{0.000}{0.000}{10.9091}{bb}{81}%
}\hskip.1em\allowbreak{} 上不可约。换元得

\[
\RB{170.334}{10.151}{3.227}{%
\RG{55.183}{4.505}{7.9701}{rm}{52}\RG{82.803}{4.505}{7.9701}{rm}{51}\RG{115.878}{4.505}{7.9701}{rm}{50}\RG{0.000}{0.000}{10.9091}{mi}{112}\RG{5.488}{0.000}{10.9091}{rm}{40}\RG{9.731}{0.000}{10.9091}{mi}{116}\RG{16.094}{0.000}{10.9091}{rm}{43}\RG{26.991}{0.000}{10.9091}{rm}{49}\RG{32.446}{0.000}{10.9091}{rm}{41}\RG{39.700}{0.000}{10.9091}{rm}{61}\RG{51.244}{0.000}{10.9091}{mi}{116}\RG{62.501}{0.000}{10.9091}{rm}{43}\RG{73.399}{0.000}{10.9091}{rm}{53}\RG{78.865}{0.000}{10.9091}{mi}{116}\RG{90.122}{0.000}{10.9091}{rm}{43}\RG{101.020}{0.000}{10.9091}{rm}{49}\RG{106.474}{0.000}{10.9091}{rm}{48}\RG{111.939}{0.000}{10.9091}{mi}{116}\RG{123.196}{0.000}{10.9091}{rm}{43}\RG{134.094}{0.000}{10.9091}{rm}{49}\RG{139.548}{0.000}{10.9091}{rm}{48}\RG{145.013}{0.000}{10.9091}{mi}{116}\RG{151.376}{0.000}{10.9091}{rm}{43}\RG{162.274}{0.000}{10.9091}{rm}{53}\RG{167.739}{0.000}{10.9091}{mi}{58}%
}
\]
其首项系数不被\allowbreak{}\hskip.1em \RB{5.398}{7.765}{0.740}{%
\RG{0.000}{-0.000}{10.9091}{rm}{53}%
}\hskip.1em\allowbreak{} 整除，其余系数均被\allowbreak{}\hskip.1em \RB{5.398}{7.765}{0.740}{%
\RG{0.000}{-0.000}{10.9091}{rm}{53}%
}\hskip.1em\allowbreak{} 整除，而常数项不被\allowbreak{}\hskip.1em \RB{10.852}{7.765}{0.740}{%
\RG{0.000}{-0.000}{10.9091}{rm}{50}\RG{5.454}{-0.000}{10.9091}{rm}{53}%
}\hskip.1em\allowbreak{} 整除。由Eisenstein判别法，\allowbreak{}\hskip.1em \RB{36.345}{8.681}{3.227}{%
\RG{0.000}{-0.000}{10.9091}{mi}{112}\RG{5.488}{-0.000}{10.9091}{rm}{40}\RG{9.730}{-0.000}{10.9091}{mi}{116}\RG{16.221}{-0.000}{10.9091}{rm}{43}\RG{27.238}{-0.000}{10.9091}{rm}{49}\RG{32.693}{-0.000}{10.9091}{rm}{41}%
}\hskip.1em\allowbreak{} 在\allowbreak{}\hskip.1em \RB{8.627}{8.180}{2.453}{%
\RG{0.000}{-0.000}{10.9091}{bb}{81}%
}\hskip.1em\allowbreak{}上不可约；平移换元保持可约性，故\allowbreak{}\hskip.1em \RB{17.322}{8.685}{3.224}{%
\RG{0.000}{0.003}{10.9091}{mi}{112}\RG{5.488}{0.003}{10.9091}{rm}{40}\RG{9.730}{0.003}{10.9091}{mi}{116}\RG{13.669}{0.003}{10.9091}{rm}{41}%
}\hskip.1em\allowbreak{} 也不可约。

若四个向量线性相关，则存在不全为零的\allowbreak{}\hskip.1em \RB{74.963}{8.185}{2.600}{%
\RG{0.000}{0.005}{10.9091}{mi}{99}\RG{4.720}{-1.630}{7.9701}{rm}{48}\RG{9.615}{0.005}{10.9091}{mi}{59}\RG{14.447}{0.005}{10.9091}{mi}{99}\RG{19.184}{-1.630}{7.9701}{rm}{49}\RG{24.077}{0.005}{10.9091}{mi}{59}\RG{28.910}{0.005}{10.9091}{mi}{99}\RG{33.646}{-1.630}{7.9701}{rm}{50}\RG{38.540}{0.005}{10.9091}{mi}{59}\RG{43.373}{0.005}{10.9091}{mi}{99}\RG{48.109}{-1.630}{7.9701}{rm}{51}\RG{56.033}{0.005}{10.9091}{sy}{50}\RG{66.336}{0.005}{10.9091}{bb}{81}%
}\hskip.1em\allowbreak{} ，使

\[
\RB{235.350}{10.080}{3.222}{%
\RG{176.082}{4.510}{7.9701}{rm}{50}\RG{207.863}{4.510}{7.9701}{rm}{51}\RG{0.000}{0.005}{10.9091}{mi}{113}\RG{5.261}{0.005}{10.9091}{rm}{40}\RG{9.504}{0.005}{10.9091}{mi}{65}\RG{17.686}{0.005}{10.9091}{rm}{41}\RG{21.928}{0.005}{10.9091}{mi}{11}\RG{36.830}{0.005}{10.9091}{rm}{61}\RG{48.328}{0.005}{10.9091}{rm}{48}\RG{53.799}{0.005}{10.9091}{mi}{59}\RG{80.449}{0.005}{10.9091}{mi}{113}\RG{85.726}{0.005}{10.9091}{rm}{40}\RG{89.967}{0.005}{10.9091}{mi}{116}\RG{93.907}{0.005}{10.9091}{rm}{41}\RG{101.161}{0.005}{10.9091}{rm}{61}\RG{112.694}{0.005}{10.9091}{mi}{99}\RG{124.733}{0.005}{10.9091}{rm}{43}\RG{135.642}{0.005}{10.9091}{mi}{99}\RG{145.256}{0.005}{10.9091}{mi}{116}\RG{151.619}{0.005}{10.9091}{rm}{43}\RG{162.528}{0.005}{10.9091}{mi}{99}\RG{172.143}{0.005}{10.9091}{mi}{116}\RG{183.400}{0.005}{10.9091}{rm}{43}\RG{194.308}{0.005}{10.9091}{mi}{99}\RG{203.923}{0.005}{10.9091}{mi}{116}\RG{215.786}{0.005}{10.9091}{sy}{54}\RG{215.786}{0.005}{10.9091}{rm}{61}\RG{227.284}{0.005}{10.9091}{rm}{48}\RG{232.756}{0.005}{10.9091}{mi}{58}\RG{28.906}{-1.630}{7.9701}{rm}{49}\RG{117.415}{-1.630}{7.9701}{rm}{48}\RG{140.363}{-1.630}{7.9701}{rm}{49}\RG{167.249}{-1.630}{7.9701}{rm}{50}\RG{199.030}{-1.630}{7.9701}{rm}{51}%
}
\]
由于deg\allowbreak{}\hskip.1em \RB{16.350}{6.386}{2.621}{%
\RG{0.000}{-0.004}{10.9091}{mi}{113}\RG{8.280}{-0.004}{10.9091}{mi}{60}%
}\hskip.1em\allowbreak{} deg\allowbreak{}\hskip.1em \RB{5.845}{5.317}{2.621}{%
\RG{0.000}{-0.004}{10.9091}{mi}{112}%
}\hskip.1em\allowbreak{} 且\allowbreak{}\hskip.1em \RB{5.845}{5.317}{2.621}{%
\RG{0.000}{-0.004}{10.9091}{mi}{112}%
}\hskip.1em\allowbreak{} 不可约，gcd\allowbreak{}\hskip.1em \RB{43.673}{8.677}{3.232}{%
\RG{0.000}{-0.004}{10.9091}{rm}{40}\RG{4.242}{-0.004}{10.9091}{mi}{112}\RG{9.729}{-0.004}{10.9091}{mi}{59}\RG{14.561}{-0.004}{10.9091}{mi}{113}\RG{19.840}{-0.004}{10.9091}{rm}{41}\RG{27.094}{-0.004}{10.9091}{rm}{61}\RG{38.602}{-0.004}{10.9091}{rm}{49}%
}\hskip.1em\allowbreak{} 。由多项式的Bézout等式，存在\allowbreak{}\hskip.1em \RB{47.794}{8.677}{3.232}{%
\RG{0.000}{-0.004}{10.9091}{mi}{117}\RG{6.240}{-0.004}{10.9091}{mi}{59}\RG{11.072}{-0.004}{10.9091}{mi}{118}\RG{19.802}{-0.004}{10.9091}{sy}{50}\RG{30.105}{-0.004}{10.9091}{bb}{81}\RG{38.590}{-0.004}{10.9091}{rm}{91}\RG{41.621}{-0.004}{10.9091}{mi}{116}\RG{45.559}{-0.004}{10.9091}{rm}{93}%
}\hskip.1em\allowbreak{} ，使\allowbreak{}\hskip.1em \RB{54.821}{7.761}{2.621}{%
\RG{0.000}{-0.004}{10.9091}{mi}{117}\RG{6.240}{-0.004}{10.9091}{mi}{112}\RG{13.766}{-0.004}{10.9091}{rm}{43}\RG{24.283}{-0.004}{10.9091}{mi}{118}\RG{29.945}{-0.004}{10.9091}{mi}{113}\RG{38.252}{-0.004}{10.9091}{rm}{61}\RG{49.750}{-0.004}{10.9091}{rm}{49}%
}\hskip.1em\allowbreak{} 。代入\allowbreak{}\hskip.1em \RB{8.365}{8.309}{0.501}{%
\RG{0.000}{-0.001}{10.9091}{mi}{65}%
}\hskip.1em\allowbreak{} 并作用于\allowbreak{}\hskip.1em \RB{11.090}{5.321}{2.138}{%
\RG{0.000}{-0.001}{10.9091}{mi}{11}\RG{6.978}{-1.638}{7.9701}{rm}{49}%
}\hskip.1em\allowbreak{} ，得到

\[
\RB{175.545}{8.684}{3.225}{%
\RG{0.000}{0.002}{10.9091}{mi}{11}\RG{14.902}{0.002}{10.9091}{rm}{61}\RG{26.417}{0.002}{10.9091}{mi}{117}\RG{32.662}{0.002}{10.9091}{rm}{40}\RG{36.903}{0.002}{10.9091}{mi}{65}\RG{45.085}{0.002}{10.9091}{rm}{41}\RG{49.328}{0.002}{10.9091}{mi}{112}\RG{54.816}{0.002}{10.9091}{rm}{40}\RG{59.059}{0.002}{10.9091}{mi}{65}\RG{67.240}{0.002}{10.9091}{rm}{41}\RG{71.482}{0.002}{10.9091}{mi}{11}\RG{85.778}{0.002}{10.9091}{rm}{43}\RG{96.687}{0.002}{10.9091}{mi}{118}\RG{102.366}{0.002}{10.9091}{rm}{40}\RG{106.609}{0.002}{10.9091}{mi}{65}\RG{114.790}{0.002}{10.9091}{rm}{41}\RG{119.032}{0.002}{10.9091}{mi}{113}\RG{124.293}{0.002}{10.9091}{rm}{40}\RG{128.535}{0.002}{10.9091}{mi}{65}\RG{136.717}{0.002}{10.9091}{rm}{41}\RG{140.960}{0.002}{10.9091}{mi}{11}\RG{155.862}{0.002}{10.9091}{rm}{61}\RG{167.360}{0.002}{10.9091}{rm}{48}\RG{172.831}{0.002}{10.9091}{mi}{59}\RG{6.978}{-1.635}{7.9701}{rm}{49}\RG{78.460}{-1.635}{7.9701}{rm}{49}\RG{147.938}{-1.635}{7.9701}{rm}{49}%
}
\]
与条件矛盾。因此四个向量线性无关，又dim\allowbreak{}\hskip.1em \RB{33.544}{9.602}{2.456}{%
\RG{0.000}{-0.003}{10.9091}{bb}{81}\RG{8.485}{3.955}{7.9701}{rm}{52}\RG{16.408}{-0.003}{10.9091}{rm}{61}\RG{27.906}{-0.003}{10.9091}{rm}{52}%
}\hskip.1em\allowbreak{} ，所以它们构成一组基。\allowbreak{}\hskip.1em \RB{8.365}{7.991}{0.503}{%
\RG{0.000}{-0.003}{10.9091}{am}{3}%
}\hskip.1em\allowbreak{}

小结：把向量递推关系写成\allowbreak{}\hskip.1em \RB{54.292}{8.677}{3.231}{%
\RG{0.000}{-0.004}{10.9091}{mi}{112}\RG{5.488}{-0.004}{10.9091}{rm}{40}\RG{9.731}{-0.004}{10.9091}{mi}{65}\RG{17.912}{-0.004}{10.9091}{rm}{41}\RG{22.154}{-0.004}{10.9091}{mi}{11}\RG{29.133}{-1.640}{7.9701}{rm}{49}\RG{37.167}{-0.004}{10.9091}{rm}{61}\RG{48.774}{-0.004}{10.9091}{rm}{48}%
}\hskip.1em\allowbreak{} 。若\allowbreak{}\hskip.1em \RB{5.845}{5.317}{2.621}{%
\RG{0.000}{-0.004}{10.9091}{mi}{112}%
}\hskip.1em\allowbreak{} 不可约，则任何次数更低的非零多项式都不能消去非零向量\allowbreak{}\hskip.1em \RB{11.090}{5.317}{2.140}{%
\RG{0.000}{-0.004}{10.9091}{mi}{11}\RG{6.978}{-1.640}{7.9701}{rm}{49}%
}\hskip.1em\allowbreak{} ，从而得到线性无关性。

% Original page 2.
\section*{2014A08平方和不变与相似变换}
题目设\allowbreak{}\hskip.1em \RB{44.070}{8.687}{3.688}{%
\RG{0.000}{0.005}{10.9091}{mi}{65}\RG{11.211}{0.005}{10.9091}{rm}{61}\RG{22.709}{0.005}{10.9091}{rm}{40}\RG{26.968}{0.005}{10.9091}{mi}{97}\RG{32.735}{-1.630}{7.9701}{mi}{105}\RG{35.820}{-1.630}{7.9701}{mi}{106}\RG{40.417}{0.005}{10.9091}{rm}{41}%
}\hskip.1em\allowbreak{} 为数域\allowbreak{}\hskip.1em \RB{8.889}{7.978}{0.500}{%
\RG{0.000}{0.005}{10.9091}{bb}{75}%
}\hskip.1em\allowbreak{} 上的\allowbreak{}\hskip.1em \RB{6.729}{5.327}{0.615}{%
\RG{0.000}{0.005}{10.9091}{mi}{110}%
}\hskip.1em\allowbreak{} 阶方阵，定义

\[
\RB{76.361}{17.507}{14.925}{%
\RG{45.417}{13.640}{7.9701}{mi}{110}\RG{40.234}{10.367}{10.9087}{ex}{88}\RG{66.084}{5.246}{7.9701}{mi}{59}\RG{68.675}{5.246}{7.9701}{rm}{50}\RG{0.000}{0.003}{10.9091}{mi}{102}\RG{6.515}{0.004}{10.9091}{rm}{40}\RG{10.757}{0.003}{10.9091}{mi}{65}\RG{18.938}{0.004}{10.9091}{rm}{41}\RG{26.193}{0.004}{10.9091}{rm}{61}\RG{60.317}{0.004}{10.9091}{mi}{97}\RG{73.766}{0.003}{10.9091}{mi}{58}\RG{66.083}{-3.036}{7.9701}{mi}{105}\RG{69.168}{-3.036}{7.9701}{mi}{106}\RG{37.726}{-12.867}{7.9701}{mi}{105}\RG{40.811}{-12.867}{7.9701}{mi}{59}\RG{43.399}{-12.867}{7.9701}{mi}{106}\RG{47.454}{-12.867}{7.9701}{rm}{61}\RG{54.151}{-12.867}{7.9701}{rm}{49}%
}
\]
设\allowbreak{}\hskip.1em \RB{8.725}{7.948}{0.502}{%
\RG{0.000}{-0.002}{10.9091}{mi}{80}%
}\hskip.1em\allowbreak{} 为数域\allowbreak{}\hskip.1em \RB{8.889}{7.970}{0.502}{%
\RG{0.000}{-0.002}{10.9091}{bb}{75}%
}\hskip.1em\allowbreak{} 上的\allowbreak{}\hskip.1em \RB{6.729}{5.319}{0.622}{%
\RG{0.000}{-0.002}{10.9091}{mi}{110}%
}\hskip.1em\allowbreak{} 阶非异阵，满足：对任意\allowbreak{}\hskip.1em \RB{54.684}{8.679}{3.229}{%
\RG{0.000}{-0.002}{10.9091}{mi}{65}\RG{11.346}{-0.002}{10.9091}{sy}{50}\RG{21.783}{-0.002}{10.9091}{mi}{77}\RG{32.367}{-1.638}{7.9701}{mi}{110}\RG{38.305}{-0.002}{10.9091}{rm}{40}\RG{42.547}{-0.002}{10.9091}{bb}{75}\RG{51.032}{-0.002}{10.9091}{rm}{41}%
}\hskip.1em\allowbreak{} ，成立\allowbreak{}\hskip.1em \RB{89.335}{9.526}{3.229}{%
\RG{0.000}{-0.002}{10.9091}{mi}{102}\RG{6.515}{-0.002}{10.9091}{rm}{40}\RG{10.758}{-0.002}{10.9091}{mi}{80}\RG{19.267}{-0.002}{10.9091}{mi}{65}\RG{27.448}{-0.002}{10.9091}{mi}{80}\RG{35.976}{3.957}{7.9701}{sy}{0}\RG{42.794}{3.956}{7.9701}{rm}{49}\RG{47.688}{-0.002}{10.9091}{rm}{41}\RG{55.084}{-0.002}{10.9091}{rm}{61}\RG{66.744}{-0.002}{10.9091}{mi}{102}\RG{73.259}{-0.002}{10.9091}{rm}{40}\RG{77.501}{-0.002}{10.9091}{mi}{65}\RG{85.683}{-0.002}{10.9091}{rm}{41}%
}\hskip.1em\allowbreak{} 。证明：存在非零常数\allowbreak{}\hskip.1em \RB{26.942}{7.974}{0.935}{%
\RG{0.000}{0.001}{10.9091}{mi}{99}\RG{7.750}{0.001}{10.9091}{sy}{50}\RG{18.054}{0.001}{10.9091}{bb}{75}%
}\hskip.1em\allowbreak{} ，使得\allowbreak{}\hskip.1em \RB{49.683}{8.728}{2.211}{%
\RG{0.000}{0.001}{10.9091}{mi}{80}\RG{8.509}{0.001}{10.9091}{mi}{80}\RG{17.038}{3.959}{7.9701}{sy}{48}\RG{23.128}{0.001}{10.9091}{rm}{61}\RG{34.642}{0.001}{10.9091}{mi}{99}\RG{39.355}{0.001}{10.9091}{mi}{73}\RG{44.159}{-1.635}{7.9701}{mi}{110}%
}\hskip.1em\allowbreak{} 。

解：由\allowbreak{}\hskip.1em \RB{34.558}{8.679}{3.229}{%
\RG{0.000}{-0.002}{10.9091}{mi}{102}\RG{6.515}{-0.002}{10.9091}{rm}{40}\RG{10.757}{-0.002}{10.9091}{mi}{65}\RG{18.938}{-0.002}{10.9091}{rm}{41}\RG{26.193}{-0.002}{10.9091}{rm}{61}%
}\hskip.1em\allowbreak{} tr\allowbreak{}\hskip.1em \RB{27.318}{8.726}{3.229}{%
\RG{0.000}{-0.002}{10.9091}{rm}{40}\RG{4.242}{-0.002}{10.9091}{mi}{65}\RG{12.423}{-0.002}{10.9091}{mi}{65}\RG{20.605}{3.957}{7.9701}{sy}{48}\RG{23.665}{-0.002}{10.9091}{rm}{41}%
}\hskip.1em\allowbreak{} ，可得极化恒等式

\[
\RB{180.483}{9.291}{3.228}{%
\RG{170.585}{4.502}{7.9701}{sy}{48}\RG{0.000}{-0.001}{10.9091}{mi}{102}\RG{6.515}{-0.001}{10.9091}{rm}{40}\RG{10.757}{-0.001}{10.9091}{mi}{65}\RG{21.363}{-0.001}{10.9091}{rm}{43}\RG{32.271}{-0.001}{10.9091}{mi}{66}\RG{41.093}{-0.001}{10.9091}{rm}{41}\RG{47.759}{-0.001}{10.9091}{sy}{0}\RG{58.668}{-0.001}{10.9091}{mi}{102}\RG{65.183}{-0.001}{10.9091}{rm}{40}\RG{69.425}{-0.001}{10.9091}{mi}{65}\RG{77.606}{-0.001}{10.9091}{rm}{41}\RG{84.273}{-0.001}{10.9091}{sy}{0}\RG{95.182}{-0.001}{10.9091}{mi}{102}\RG{101.697}{-0.001}{10.9091}{rm}{40}\RG{105.938}{-0.001}{10.9091}{mi}{66}\RG{114.761}{-0.001}{10.9091}{rm}{41}\RG{122.015}{-0.001}{10.9091}{rm}{61}\RG{133.523}{-0.001}{10.9091}{rm}{50}\RT{140.820}{-0.001}{10.9091}{t}\RT{145.063}{-0.001}{10.9091}{r}\RG{149.339}{-0.001}{10.9091}{rm}{40}\RG{153.582}{-0.001}{10.9091}{mi}{65}\RG{161.764}{-0.001}{10.9091}{mi}{66}\RG{173.645}{-0.001}{10.9091}{rm}{41}\RG{177.888}{-0.001}{10.9091}{mi}{58}%
}
\]
将假设分别用于\allowbreak{}\hskip.1em \RB{56.932}{8.309}{2.606}{%
\RG{0.000}{-0.001}{10.9091}{mi}{65}\RG{10.606}{-0.001}{10.9091}{rm}{43}\RG{21.514}{-0.001}{10.9091}{mi}{66}\RG{30.328}{-0.001}{10.9091}{mi}{59}\RG{35.160}{-0.001}{10.9091}{mi}{65}\RG{43.342}{-0.001}{10.9091}{mi}{59}\RG{48.185}{-0.001}{10.9091}{mi}{66}%
}\hskip.1em\allowbreak{} （数域的特征为\allowbreak{}\hskip.1em \RB{5.518}{7.764}{0.741}{%
\RG{0.000}{-0.001}{10.9091}{rm}{48}%
}\hskip.1em\allowbreak{} ，故可除以\allowbreak{}\hskip.1em \RB{5.398}{7.764}{0.501}{%
\RG{0.000}{-0.001}{10.9091}{rm}{50}%
}\hskip.1em\allowbreak{} ），得到

\[
\RB{266.234}{10.074}{4.308}{%
\RG{8.519}{8.835}{10.9087}{ex}{0}\RG{108.050}{8.835}{10.9087}{ex}{1}\RG{42.980}{4.504}{7.9701}{sy}{0}\RG{49.798}{4.504}{7.9701}{rm}{49}\RG{89.036}{4.504}{7.9701}{sy}{0}\RG{95.854}{4.504}{7.9701}{rm}{49}\RG{104.989}{4.504}{7.9701}{sy}{48}\RG{157.360}{4.504}{7.9701}{sy}{48}\RT{0.000}{-0.000}{10.9091}{t}\RT{4.243}{-0.000}{10.9091}{r}\RG{13.519}{-0.000}{10.9091}{rm}{40}\RG{17.761}{-0.000}{10.9091}{mi}{80}\RG{26.270}{-0.000}{10.9091}{mi}{65}\RG{34.452}{-0.000}{10.9091}{mi}{80}\RG{54.692}{-0.000}{10.9091}{rm}{41}\RG{58.925}{-0.000}{10.9091}{rm}{40}\RG{63.177}{-0.000}{10.9091}{mi}{80}\RG{71.686}{-0.000}{10.9091}{mi}{66}\RG{80.511}{-0.000}{10.9091}{mi}{80}\RG{100.747}{-0.000}{10.9091}{rm}{41}\RG{116.080}{-0.000}{10.9091}{rm}{61}\RT{127.595}{-0.000}{10.9091}{t}\RT{131.838}{-0.000}{10.9091}{r}\RG{136.115}{-0.000}{10.9091}{rm}{40}\RG{140.356}{-0.000}{10.9091}{mi}{65}\RG{148.538}{-0.000}{10.9091}{mi}{66}\RG{160.420}{-0.000}{10.9091}{rm}{41}\RG{186.459}{-0.000}{10.9091}{rm}{40}\RG{190.722}{-0.000}{10.9091}{mi}{65}\RG{198.904}{-0.000}{10.9091}{mi}{59}\RG{203.736}{-0.000}{10.9091}{mi}{66}\RG{215.604}{-0.000}{10.9091}{sy}{50}\RG{225.907}{-0.000}{10.9091}{mi}{77}\RG{242.428}{-0.000}{10.9091}{rm}{40}\RG{246.669}{-0.000}{10.9091}{bb}{75}\RG{255.154}{-0.000}{10.9091}{rm}{41}\RG{259.387}{-0.000}{10.9091}{rm}{41}\RG{263.639}{-0.000}{10.9091}{mi}{58}\RG{236.490}{-1.637}{7.9701}{mi}{110}%
}
\]
令\allowbreak{}\hskip.1em \RB{42.167}{8.727}{0.740}{%
\RG{0.000}{-0.000}{10.9091}{mi}{83}\RG{10.348}{-0.000}{10.9091}{rm}{61}\RG{21.863}{-0.000}{10.9091}{mi}{80}\RG{30.382}{3.958}{7.9701}{sy}{48}\RG{33.442}{-0.000}{10.9091}{mi}{80}%
}\hskip.1em\allowbreak{} 。注意\allowbreak{}\hskip.1em \RB{84.798}{9.529}{3.227}{%
\RG{0.000}{-0.000}{10.9091}{mi}{80}\RG{8.519}{3.958}{7.9701}{sy}{0}\RG{15.337}{3.958}{7.9701}{rm}{49}\RG{20.230}{-0.000}{10.9091}{rm}{40}\RG{24.472}{-0.000}{10.9091}{mi}{80}\RG{32.991}{3.958}{7.9701}{sy}{0}\RG{39.809}{3.958}{7.9701}{rm}{49}\RG{44.703}{-0.000}{10.9091}{rm}{41}\RG{48.946}{3.958}{7.9701}{sy}{48}\RG{55.036}{-0.000}{10.9091}{rm}{61}\RG{66.551}{-0.000}{10.9091}{mi}{83}\RG{73.868}{3.958}{7.9701}{sy}{0}\RG{80.687}{3.958}{7.9701}{rm}{49}%
}\hskip.1em\allowbreak{} ，利用迹的循环性质，上式化为

tr\allowbreak{}\hskip.1em \RB{67.806}{10.075}{4.307}{%
\RG{0.000}{8.836}{10.9087}{ex}{0}\RG{5.000}{-0.000}{10.9091}{mi}{65}\RG{13.181}{-0.000}{10.9091}{mi}{83}\RG{20.500}{4.505}{7.9701}{sy}{0}\RG{27.317}{4.505}{7.9701}{rm}{49}\RG{32.212}{-0.000}{10.9091}{mi}{66}\RG{41.033}{4.505}{7.9701}{sy}{48}\RG{44.093}{-0.000}{10.9091}{mi}{83}\RG{51.412}{8.836}{10.9087}{ex}{1}\RG{59.441}{-0.000}{10.9091}{rm}{61}%
}\hskip.1em\allowbreak{} tr\allowbreak{}\hskip.1em \RB{31.263}{9.273}{3.227}{%
\RG{0.000}{-0.000}{10.9091}{rm}{40}\RG{4.243}{-0.000}{10.9091}{mi}{65}\RG{12.424}{-0.000}{10.9091}{mi}{66}\RG{21.246}{4.505}{7.9701}{sy}{48}\RG{24.306}{-0.000}{10.9091}{rm}{41}\RG{28.549}{-0.000}{10.9091}{mi}{59}%
}\hskip.1em\allowbreak{} 即tr\allowbreak{}\hskip.1em \RB{112.704}{10.075}{4.307}{%
\RG{0.000}{8.836}{10.9087}{ex}{0}\RG{5.000}{-0.000}{10.9091}{mi}{65}\RG{13.181}{-0.000}{10.9091}{rm}{40}\RG{17.423}{-0.000}{10.9091}{mi}{83}\RG{24.741}{4.505}{7.9701}{sy}{0}\RG{31.559}{4.505}{7.9701}{rm}{49}\RG{36.452}{-0.000}{10.9091}{mi}{66}\RG{45.274}{4.505}{7.9701}{sy}{48}\RG{48.335}{-0.000}{10.9091}{mi}{83}\RG{58.077}{-0.000}{10.9091}{sy}{0}\RG{68.986}{-0.000}{10.9091}{mi}{66}\RG{77.807}{4.505}{7.9701}{sy}{48}\RG{80.868}{-0.000}{10.9091}{rm}{41}\RG{85.110}{8.836}{10.9087}{ex}{1}\RG{93.140}{-0.000}{10.9091}{rm}{61}\RG{104.637}{-0.000}{10.9091}{rm}{48}\RG{110.109}{-0.000}{10.9091}{mi}{58}%
}\hskip.1em\allowbreak{}

对固定的\allowbreak{}\hskip.1em \RB{8.747}{7.952}{0.500}{%
\RG{0.000}{0.001}{10.9091}{mi}{66}%
}\hskip.1em\allowbreak{} ，记\allowbreak{}\hskip.1em \RB{86.289}{9.530}{0.739}{%
\RG{0.000}{0.001}{10.9091}{mi}{67}\RG{11.608}{0.001}{10.9091}{rm}{61}\RG{23.122}{0.001}{10.9091}{mi}{83}\RG{30.440}{3.959}{7.9701}{sy}{0}\RG{37.258}{3.959}{7.9701}{rm}{49}\RG{42.152}{0.001}{10.9091}{mi}{66}\RG{50.974}{3.959}{7.9701}{sy}{48}\RG{54.034}{0.001}{10.9091}{mi}{83}\RG{63.776}{0.001}{10.9091}{sy}{0}\RG{74.685}{0.001}{10.9091}{mi}{66}\RG{83.506}{3.959}{7.9701}{sy}{48}%
}\hskip.1em\allowbreak{} 。取\allowbreak{}\hskip.1em \RB{37.717}{8.312}{3.694}{%
\RG{0.000}{0.001}{10.9091}{mi}{65}\RG{11.212}{0.001}{10.9091}{rm}{61}\RG{22.727}{0.001}{10.9091}{mi}{69}\RG{30.779}{-1.636}{7.9701}{mi}{105}\RG{33.864}{-1.636}{7.9701}{mi}{106}%
}\hskip.1em\allowbreak{} ，其中\allowbreak{}\hskip.1em \RB{14.990}{7.919}{3.694}{%
\RG{0.000}{0.001}{10.9091}{mi}{69}\RG{8.052}{-1.636}{7.9701}{mi}{105}\RG{11.137}{-1.636}{7.9701}{mi}{106}%
}\hskip.1em\allowbreak{} 是第\allowbreak{}\hskip.1em \RB{21.619}{8.683}{3.226}{%
\RG{0.000}{0.001}{10.9091}{rm}{40}\RG{4.243}{0.001}{10.9091}{mi}{105}\RG{7.996}{0.001}{10.9091}{mi}{59}\RG{12.828}{0.001}{10.9091}{mi}{106}\RG{17.966}{0.001}{10.9091}{rm}{41}%
}\hskip.1em\allowbreak{} 项为\allowbreak{}\hskip.1em \RB{5.071}{7.766}{0.500}{%
\RG{0.000}{0.001}{10.9091}{rm}{49}%
}\hskip.1em\allowbreak{} 、其余项为\allowbreak{}\hskip.1em \RB{5.518}{7.766}{0.739}{%
\RG{0.000}{0.001}{10.9091}{rm}{48}%
}\hskip.1em\allowbreak{} 的矩阵，则

\[
\RB{179.052}{9.293}{3.693}{%
\RG{0.000}{0.001}{10.9091}{rm}{48}\RG{8.476}{0.001}{10.9091}{rm}{61}\RT{19.999}{0.001}{10.9091}{t}\RT{24.243}{0.001}{10.9091}{r}\RG{28.519}{0.001}{10.9091}{rm}{40}\RG{32.761}{0.001}{10.9091}{mi}{69}\RG{48.497}{0.001}{10.9091}{mi}{67}\RG{57.074}{0.001}{10.9091}{rm}{41}\RG{64.328}{0.001}{10.9091}{rm}{61}\RG{75.861}{0.001}{10.9091}{mi}{67}\RG{113.158}{0.001}{10.9091}{rm}{40}\RG{117.390}{0.001}{10.9091}{rm}{49}\RG{125.884}{0.001}{10.9091}{sy}{20}\RG{137.399}{0.001}{10.9091}{mi}{105}\RG{141.151}{0.001}{10.9091}{mi}{59}\RG{145.984}{0.001}{10.9091}{mi}{106}\RG{154.153}{0.001}{10.9091}{sy}{20}\RG{165.667}{0.001}{10.9091}{mi}{110}\RG{172.215}{0.001}{10.9091}{rm}{41}\RG{176.457}{0.001}{10.9091}{mi}{58}\RG{40.814}{-1.635}{7.9701}{mi}{105}\RG{43.899}{-1.635}{7.9701}{mi}{106}\RG{83.657}{-1.635}{7.9701}{mi}{106}\RG{87.697}{-1.635}{7.9701}{mi}{105}%
}
\]
所以\allowbreak{}\hskip.1em \RB{28.622}{8.192}{0.739}{%
\RG{0.000}{0.001}{10.9091}{mi}{67}\RG{11.607}{0.001}{10.9091}{rm}{61}\RG{23.104}{0.001}{10.9091}{rm}{48}%
}\hskip.1em\allowbreak{} ，即\allowbreak{}\hskip.1em \RB{64.380}{9.531}{0.739}{%
\RG{0.000}{0.001}{10.9091}{mi}{83}\RG{7.319}{3.960}{7.9701}{sy}{0}\RG{14.136}{3.960}{7.9701}{rm}{49}\RG{19.030}{0.001}{10.9091}{mi}{66}\RG{27.851}{3.960}{7.9701}{sy}{48}\RG{30.912}{0.001}{10.9091}{mi}{83}\RG{41.260}{0.001}{10.9091}{rm}{61}\RG{52.775}{0.001}{10.9091}{mi}{66}\RG{61.597}{3.960}{7.9701}{sy}{48}%
}\hskip.1em\allowbreak{} ，从而\allowbreak{}\hskip.1em \RB{52.667}{8.729}{0.739}{%
\RG{0.000}{0.001}{10.9091}{mi}{66}\RG{8.822}{3.960}{7.9701}{sy}{48}\RG{11.883}{0.001}{10.9091}{mi}{83}\RG{22.230}{0.001}{10.9091}{rm}{61}\RG{33.746}{0.001}{10.9091}{mi}{83}\RG{41.054}{0.001}{10.9091}{mi}{66}\RG{49.884}{3.960}{7.9701}{sy}{48}%
}\hskip.1em\allowbreak{} 。由于\allowbreak{}\hskip.1em \RB{11.605}{8.729}{0.500}{%
\RG{0.000}{0.001}{10.9091}{mi}{66}\RG{8.822}{3.960}{7.9701}{sy}{48}%
}\hskip.1em\allowbreak{} 可取任意矩阵，\allowbreak{}\hskip.1em \RB{7.536}{8.192}{0.739}{%
\RG{0.000}{0.001}{10.9091}{mi}{83}%
}\hskip.1em\allowbreak{} 与所有\allowbreak{}\hskip.1em \RB{6.729}{5.323}{0.619}{%
\RG{0.000}{0.001}{10.9091}{mi}{110}%
}\hskip.1em\allowbreak{} 阶矩阵交换。

下面说明这迫使\allowbreak{}\hskip.1em \RB{7.536}{8.194}{0.737}{%
\RG{0.000}{0.003}{10.9091}{mi}{83}%
}\hskip.1em\allowbreak{} 为数量矩阵。由\allowbreak{}\hskip.1em \RB{59.700}{8.194}{2.209}{%
\RG{0.000}{0.003}{10.9091}{mi}{83}\RG{7.309}{0.003}{10.9091}{mi}{69}\RG{15.371}{-1.633}{7.9701}{mi}{105}\RG{18.456}{-1.633}{7.9701}{mi}{105}\RG{25.498}{0.003}{10.9091}{rm}{61}\RG{37.393}{0.003}{10.9091}{mi}{69}\RG{45.445}{-1.633}{7.9701}{mi}{105}\RG{48.530}{-1.633}{7.9701}{mi}{105}\RG{52.164}{0.003}{10.9091}{mi}{83}%
}\hskip.1em\allowbreak{} ，\allowbreak{}\hskip.1em \RB{7.536}{8.194}{0.737}{%
\RG{0.000}{0.003}{10.9091}{mi}{83}%
}\hskip.1em\allowbreak{} 的第\allowbreak{}\hskip.1em \RB{3.696}{7.714}{0.617}{%
\RG{0.000}{0.003}{10.9091}{mi}{105}%
}\hskip.1em\allowbreak{} 行、第\allowbreak{}\hskip.1em \RB{3.696}{7.714}{0.617}{%
\RG{0.000}{0.003}{10.9091}{mi}{105}%
}\hskip.1em\allowbreak{} 列的非对角元均为零，故\allowbreak{}\hskip.1em \RB{7.536}{8.194}{0.737}{%
\RG{0.000}{0.003}{10.9091}{mi}{83}%
}\hskip.1em\allowbreak{} 是对角矩阵，写成

\[
\RB{98.248}{9.291}{3.228}{%
\RG{0.000}{-0.001}{10.9091}{mi}{83}\RG{10.348}{-0.001}{10.9091}{rm}{61}\RT{21.863}{-0.001}{10.9091}{d}\RT{27.928}{-0.001}{10.9091}{i}\RT{30.961}{-0.001}{10.9091}{a}\RT{36.415}{-0.001}{10.9091}{g}\RG{41.869}{-0.001}{10.9091}{rm}{40}\RG{46.111}{-0.001}{10.9091}{mi}{115}\RG{56.119}{-0.001}{10.9091}{mi}{59}\RG{60.951}{-0.001}{10.9091}{mi}{58}\RG{65.795}{-0.001}{10.9091}{mi}{58}\RG{70.627}{-0.001}{10.9091}{mi}{58}\RG{75.471}{-0.001}{10.9091}{mi}{59}\RG{80.314}{-0.001}{10.9091}{mi}{115}\RG{91.412}{-0.001}{10.9091}{rm}{41}\RG{95.654}{-0.001}{10.9091}{mi}{58}\RG{51.224}{-1.637}{7.9701}{rm}{49}\RG{85.474}{-1.637}{7.9701}{mi}{110}%
}
\]
再由\allowbreak{}\hskip.1em \RB{60.869}{8.195}{3.689}{%
\RG{0.000}{0.004}{10.9091}{mi}{83}\RG{7.309}{0.004}{10.9091}{mi}{69}\RG{15.370}{-1.632}{7.9701}{mi}{105}\RG{18.455}{-1.632}{7.9701}{mi}{106}\RG{26.083}{0.004}{10.9091}{rm}{61}\RG{37.598}{0.004}{10.9091}{mi}{69}\RG{45.651}{-1.632}{7.9701}{mi}{105}\RG{48.736}{-1.632}{7.9701}{mi}{106}\RG{53.333}{0.004}{10.9091}{mi}{83}%
}\hskip.1em\allowbreak{} ，得到\allowbreak{}\hskip.1em \RB{63.725}{7.922}{3.689}{%
\RG{0.000}{0.004}{10.9091}{mi}{115}\RG{5.113}{-1.632}{7.9701}{mi}{105}\RG{8.745}{0.004}{10.9091}{mi}{69}\RG{16.798}{-1.632}{7.9701}{mi}{105}\RG{19.883}{-1.632}{7.9701}{mi}{106}\RG{27.511}{0.004}{10.9091}{rm}{61}\RG{39.025}{0.004}{10.9091}{mi}{115}\RG{44.139}{-1.632}{7.9701}{mi}{106}\RG{48.736}{0.004}{10.9091}{mi}{69}\RG{56.788}{-1.632}{7.9701}{mi}{105}\RG{59.873}{-1.632}{7.9701}{mi}{106}%
}\hskip.1em\allowbreak{} ，即\allowbreak{}\hskip.1em \RB{32.256}{5.326}{3.689}{%
\RG{0.000}{0.004}{10.9091}{mi}{115}\RG{5.113}{-1.632}{7.9701}{mi}{105}\RG{11.775}{0.004}{10.9091}{rm}{61}\RG{23.290}{0.004}{10.9091}{mi}{115}\RG{28.403}{-1.632}{7.9701}{mi}{106}%
}\hskip.1em\allowbreak{} 。因此\allowbreak{}\hskip.1em \RB{36.904}{8.195}{2.208}{%
\RG{0.000}{0.004}{10.9091}{mi}{83}\RG{10.348}{0.004}{10.9091}{rm}{61}\RG{21.863}{0.004}{10.9091}{mi}{99}\RG{26.576}{0.004}{10.9091}{mi}{73}\RG{31.379}{-1.632}{7.9701}{mi}{110}%
}\hskip.1em\allowbreak{} ；又因\allowbreak{}\hskip.1em \RB{7.536}{8.195}{0.736}{%
\RG{0.000}{0.004}{10.9091}{mi}{83}%
}\hskip.1em\allowbreak{} 可逆，\allowbreak{}\hskip.1em \RB{24.767}{8.315}{2.852}{%
\RG{0.000}{0.004}{10.9091}{mi}{99}\RG{7.751}{0.004}{10.9091}{sy}{54}\RG{7.751}{0.004}{10.9091}{rm}{61}\RG{19.249}{0.004}{10.9091}{rm}{48}%
}\hskip.1em\allowbreak{} 。最后

\[
\RB{223.082}{10.074}{2.213}{%
\RG{8.519}{4.504}{7.9701}{sy}{48}\RG{89.516}{4.504}{7.9701}{sy}{48}\RG{120.361}{4.504}{7.9701}{sy}{0}\RG{127.179}{4.504}{7.9701}{rm}{49}\RG{183.719}{4.504}{7.9701}{sy}{48}\RG{0.000}{-0.000}{10.9091}{mi}{80}\RG{11.579}{-0.000}{10.9091}{mi}{80}\RG{23.128}{-0.000}{10.9091}{rm}{61}\RG{34.642}{-0.000}{10.9091}{mi}{99}\RG{39.355}{-0.000}{10.9091}{mi}{73}\RG{56.759}{-0.000}{10.9091}{rm}{61}\RG{63.426}{-0.000}{10.9091}{sy}{41}\RG{80.998}{-0.000}{10.9091}{mi}{80}\RG{95.606}{-0.000}{10.9091}{rm}{61}\RG{107.121}{-0.000}{10.9091}{mi}{99}\RG{111.834}{-0.000}{10.9091}{mi}{80}\RG{138.736}{-0.000}{10.9091}{rm}{61}\RG{145.402}{-0.000}{10.9091}{sy}{41}\RG{166.683}{-0.000}{10.9091}{mi}{80}\RG{175.191}{-0.000}{10.9091}{mi}{80}\RG{189.810}{-0.000}{10.9091}{rm}{61}\RG{201.325}{-0.000}{10.9091}{mi}{99}\RG{206.037}{-0.000}{10.9091}{mi}{73}\RG{220.488}{0.000}{10.9091}{mi}{58}\RG{44.159}{-1.637}{7.9701}{mi}{110}\RG{210.841}{-1.637}{7.9701}{mi}{110}%
}
\]
小结：先用极化把“平方和不变”转为双线性形式不变，再用迹配对的非退化性推出交换关系。与所有矩阵都交换的矩阵只能是数量矩阵。

% Original page 3.
\section*{2014A09从矩形乘积AB求BA}
题目设\allowbreak{}\hskip.1em \RB{21.761}{8.306}{2.610}{%
\RG{0.000}{-0.004}{10.9091}{mi}{65}\RG{8.181}{-0.004}{10.9091}{mi}{59}\RG{13.014}{-0.004}{10.9091}{mi}{66}%
}\hskip.1em\allowbreak{} 分别是\allowbreak{}\hskip.1em \RB{55.180}{7.761}{2.610}{%
\RG{0.000}{-0.004}{10.9091}{rm}{51}\RG{7.878}{-0.004}{10.9091}{sy}{2}\RG{18.787}{-0.004}{10.9091}{rm}{50}\RG{24.241}{-0.004}{10.9091}{mi}{59}\RG{30.907}{-0.004}{10.9091}{rm}{50}\RG{38.787}{-0.004}{10.9091}{sy}{2}\RG{49.695}{-0.004}{10.9091}{rm}{51}%
}\hskip.1em\allowbreak{} 矩阵，满足

\[
\RB{110.200}{11.045}{35.749}{%
\RG{31.548}{10.120}{10.9087}{ex}{48}\RG{96.243}{10.120}{10.9087}{ex}{49}\RG{45.335}{0.005}{10.9091}{rm}{56}\RG{65.930}{0.005}{10.9091}{rm}{50}\RG{82.304}{0.005}{10.9091}{sy}{0}\RG{90.788}{0.005}{10.9091}{rm}{50}\RG{31.548}{-9.079}{10.9087}{ex}{66}\RG{96.243}{-9.079}{10.9087}{ex}{67}\RG{0.000}{-15.080}{10.9091}{mi}{65}\RG{8.181}{-15.080}{10.9091}{mi}{66}\RG{20.033}{-15.080}{10.9091}{rm}{61}\RG{107.605}{-15.080}{10.9091}{mi}{58}\RG{45.335}{-15.441}{10.9091}{rm}{50}\RG{65.930}{-15.441}{10.9091}{rm}{53}\RG{86.537}{-15.441}{10.9091}{rm}{52}\RG{31.548}{-16.061}{10.9087}{ex}{64}\RG{96.243}{-16.061}{10.9087}{ex}{65}\RG{41.093}{-30.888}{10.9091}{sy}{0}\RG{49.578}{-30.888}{10.9091}{rm}{50}\RG{65.930}{-30.888}{10.9091}{rm}{52}\RG{86.536}{-30.888}{10.9091}{rm}{53}%
}
\]
试求\allowbreak{}\hskip.1em \RB{17.179}{8.309}{0.501}{%
\RG{0.000}{-0.001}{10.9091}{mi}{66}\RG{8.814}{-0.001}{10.9091}{mi}{65}%
}\hskip.1em\allowbreak{} 。

解：记\allowbreak{}\hskip.1em \RB{43.246}{8.307}{0.503}{%
\RG{0.000}{-0.003}{10.9091}{mi}{77}\RG{14.803}{-0.003}{10.9091}{rm}{61}\RG{26.317}{-0.003}{10.9091}{mi}{65}\RG{34.499}{-0.003}{10.9091}{mi}{66}%
}\hskip.1em\allowbreak{} 。直接计算得

\[
\RB{155.567}{11.043}{35.751}{%
\RG{31.211}{10.117}{10.9087}{ex}{48}\RG{112.268}{10.117}{10.9087}{ex}{49}\RG{44.999}{0.003}{10.9091}{rm}{55}\RG{50.453}{0.003}{10.9091}{rm}{50}\RG{71.048}{0.003}{10.9091}{rm}{49}\RG{76.503}{0.003}{10.9091}{rm}{56}\RG{92.876}{0.003}{10.9091}{sy}{0}\RG{101.360}{0.003}{10.9091}{rm}{49}\RG{106.815}{0.003}{10.9091}{rm}{56}\RG{31.211}{-9.082}{10.9087}{ex}{66}\RG{112.268}{-9.082}{10.9087}{ex}{67}\RG{11.772}{-10.577}{7.9701}{rm}{50}\RG{0.000}{-15.082}{10.9091}{mi}{77}\RG{19.696}{-15.082}{10.9091}{rm}{61}\RG{124.844}{-15.082}{10.9091}{rm}{61}\RG{136.342}{-15.082}{10.9091}{rm}{57}\RG{141.813}{-15.082}{10.9091}{mi}{77}\RG{152.972}{-15.082}{10.9091}{mi}{58}\RG{44.998}{-15.443}{10.9091}{rm}{49}\RG{50.453}{-15.443}{10.9091}{rm}{56}\RG{71.048}{-15.443}{10.9091}{rm}{52}\RG{76.503}{-15.443}{10.9091}{rm}{53}\RG{97.109}{-15.443}{10.9091}{rm}{51}\RG{102.563}{-15.443}{10.9091}{rm}{54}\RG{31.211}{-16.063}{10.9087}{ex}{64}\RG{112.268}{-16.063}{10.9087}{ex}{65}\RG{40.756}{-30.890}{10.9091}{sy}{0}\RG{49.240}{-30.890}{10.9091}{rm}{49}\RG{54.695}{-30.890}{10.9091}{rm}{56}\RG{71.058}{-30.890}{10.9091}{rm}{51}\RG{76.512}{-30.890}{10.9091}{rm}{54}\RG{97.108}{-30.890}{10.9091}{rm}{52}\RG{102.562}{-30.890}{10.9091}{rm}{53}%
}
\]
又rank\allowbreak{}\hskip.1em \RB{31.714}{7.953}{1.992}{%
\RG{0.000}{0.002}{10.9091}{mi}{77}\RG{14.802}{0.002}{10.9091}{sy}{20}\RG{26.316}{0.002}{10.9091}{rm}{50}%
}\hskip.1em\allowbreak{} ，而其左上角二阶子式为\allowbreak{}\hskip.1em \RB{96.407}{8.313}{2.854}{%
\RG{0.000}{0.002}{10.9091}{rm}{56}\RG{7.879}{0.002}{10.9091}{sy}{1}\RG{13.333}{0.002}{10.9091}{rm}{53}\RG{21.212}{0.002}{10.9091}{sy}{0}\RG{32.121}{0.002}{10.9091}{rm}{50}\RG{39.999}{0.002}{10.9091}{sy}{1}\RG{45.453}{0.002}{10.9091}{rm}{50}\RG{53.929}{0.002}{10.9091}{rm}{61}\RG{65.438}{0.002}{10.9091}{rm}{51}\RG{70.892}{0.002}{10.9091}{rm}{54}\RG{79.391}{0.002}{10.9091}{sy}{54}\RG{79.391}{0.002}{10.9091}{rm}{61}\RG{90.889}{0.002}{10.9091}{rm}{48}%
}\hskip.1em\allowbreak{} ，故rank\allowbreak{}\hskip.1em \RB{31.697}{7.953}{0.500}{%
\RG{0.000}{0.002}{10.9091}{mi}{77}\RG{14.802}{0.002}{10.9091}{rm}{61}\RG{26.299}{0.002}{10.9091}{rm}{50}%
}\hskip.1em\allowbreak{} 。于是

\[
\RB{201.993}{9.295}{3.225}{%
\RG{0.000}{0.002}{10.9091}{rm}{50}\RG{8.476}{0.002}{10.9091}{rm}{61}\RT{19.999}{0.002}{10.9091}{r}\RT{24.276}{0.002}{10.9091}{a}\RT{29.730}{0.002}{10.9091}{n}\RT{35.795}{0.002}{10.9091}{k}\RG{41.555}{0.002}{10.9091}{rm}{40}\RG{45.797}{0.002}{10.9091}{mi}{65}\RG{53.978}{0.002}{10.9091}{mi}{66}\RG{62.800}{0.002}{10.9091}{rm}{41}\RG{70.072}{0.002}{10.9091}{sy}{20}\RT{81.588}{0.002}{10.9091}{m}\RT{90.675}{0.002}{10.9091}{i}\RT{93.707}{0.002}{10.9091}{n}\RG{99.772}{0.002}{10.9091}{sy}{102}\RT{105.227}{0.002}{10.9091}{r}\RT{109.503}{0.002}{10.9091}{a}\RT{114.957}{0.002}{10.9091}{n}\RT{121.022}{0.002}{10.9091}{k}\RG{128.600}{0.002}{10.9091}{mi}{65}\RG{136.781}{0.002}{10.9091}{mi}{59}\RT{141.630}{0.002}{10.9091}{r}\RT{145.906}{0.002}{10.9091}{a}\RT{151.360}{0.002}{10.9091}{n}\RT{157.426}{0.002}{10.9091}{k}\RG{165.004}{0.002}{10.9091}{mi}{66}\RG{173.825}{0.002}{10.9091}{sy}{103}\RG{182.301}{0.002}{10.9091}{sy}{20}\RG{193.824}{0.002}{10.9091}{rm}{50}\RG{199.278}{0.002}{10.9091}{mi}{59}%
}
\]
即\allowbreak{}\hskip.1em \RB{8.365}{8.314}{0.500}{%
\RG{0.000}{0.003}{10.9091}{mi}{65}%
}\hskip.1em\allowbreak{} 列满秩、\allowbreak{}\hskip.1em \RB{8.747}{7.954}{0.500}{%
\RG{0.000}{0.003}{10.9091}{mi}{66}%
}\hskip.1em\allowbreak{} 行满秩。因此存在左逆\allowbreak{}\hskip.1em \RB{7.514}{7.954}{0.500}{%
\RG{0.000}{0.003}{10.9091}{mi}{76}%
}\hskip.1em\allowbreak{} 和右逆\allowbreak{}\hskip.1em \RB{8.736}{7.954}{0.737}{%
\RG{0.000}{0.003}{10.9091}{mi}{82}%
}\hskip.1em\allowbreak{} ，满足\allowbreak{}\hskip.1em \RB{39.302}{8.314}{2.134}{%
\RG{0.000}{0.003}{10.9091}{mi}{76}\RG{7.418}{0.003}{10.9091}{mi}{65}\RG{18.636}{0.003}{10.9091}{rm}{61}\RG{30.151}{0.003}{10.9091}{mi}{73}\RG{34.946}{-1.634}{7.9701}{rm}{50}%
}\hskip.1em\allowbreak{} 、\allowbreak{}\hskip.1em \RB{40.885}{7.954}{2.134}{%
\RG{0.000}{0.003}{10.9091}{mi}{66}\RG{8.814}{0.003}{10.9091}{mi}{82}\RG{20.220}{0.003}{10.9091}{rm}{61}\RG{31.734}{0.003}{10.9091}{mi}{73}\RG{36.529}{-1.634}{7.9701}{rm}{50}%
}\hskip.1em\allowbreak{} 。

由\allowbreak{}\hskip.1em \RB{48.553}{9.534}{0.736}{%
\RG{0.000}{0.004}{10.9091}{mi}{77}\RG{11.772}{3.964}{7.9701}{rm}{50}\RG{19.696}{0.004}{10.9091}{rm}{61}\RG{31.194}{0.004}{10.9091}{rm}{57}\RG{36.665}{0.004}{10.9091}{mi}{77}%
}\hskip.1em\allowbreak{} ，有\allowbreak{}\hskip.1em \RB{91.013}{8.686}{3.223}{%
\RG{0.000}{0.004}{10.9091}{mi}{65}\RG{8.182}{0.004}{10.9091}{rm}{40}\RG{12.423}{0.004}{10.9091}{mi}{66}\RG{21.238}{0.004}{10.9091}{mi}{65}\RG{31.851}{0.004}{10.9091}{sy}{0}\RG{42.759}{0.004}{10.9091}{rm}{57}\RG{48.213}{0.004}{10.9091}{mi}{73}\RG{53.009}{-1.631}{7.9701}{rm}{50}\RG{57.903}{0.004}{10.9091}{rm}{41}\RG{62.146}{0.004}{10.9091}{mi}{66}\RG{73.997}{0.004}{10.9091}{rm}{61}\RG{85.495}{0.004}{10.9091}{rm}{48}%
}\hskip.1em\allowbreak{} 。左乘\allowbreak{}\hskip.1em \RB{7.514}{7.955}{0.500}{%
\RG{0.000}{0.004}{10.9091}{mi}{76}%
}\hskip.1em\allowbreak{} 、右乘\allowbreak{}\hskip.1em \RB{8.736}{7.955}{0.736}{%
\RG{0.000}{0.004}{10.9091}{mi}{82}%
}\hskip.1em\allowbreak{} ，得到

\[
\RB{172.795}{19.593}{14.122}{%
\RG{127.402}{18.657}{10.9087}{ex}{32}\RG{157.855}{18.657}{10.9087}{ex}{33}\RG{136.038}{7.364}{10.9091}{rm}{57}\RG{152.390}{7.364}{10.9091}{rm}{48}\RG{0.000}{0.003}{10.9091}{mi}{66}\RG{8.814}{0.003}{10.9091}{mi}{65}\RG{19.428}{0.003}{10.9091}{sy}{0}\RG{30.337}{0.003}{10.9091}{rm}{57}\RG{35.791}{0.003}{10.9091}{mi}{73}\RG{48.510}{0.003}{10.9091}{rm}{61}\RG{60.008}{0.003}{10.9091}{rm}{48}\RG{65.479}{0.003}{10.9091}{mi}{59}\RG{95.837}{0.003}{10.9091}{mi}{66}\RG{104.652}{0.003}{10.9091}{mi}{65}\RG{115.887}{0.003}{10.9091}{rm}{61}\RG{170.201}{0.003}{10.9091}{mi}{58}\RG{40.586}{-1.634}{7.9701}{rm}{50}\RG{136.038}{-8.083}{10.9091}{rm}{48}\RG{152.390}{-8.083}{10.9091}{rm}{57}%
}
\]
小结：先寻找\allowbreak{}\hskip.1em \RB{16.928}{8.308}{0.502}{%
\RG{0.000}{-0.002}{10.9091}{mi}{65}\RG{8.181}{-0.002}{10.9091}{mi}{66}%
}\hskip.1em\allowbreak{} 满足的低次多项式关系，再利用满秩矩形矩阵的左右逆，将关系传给\allowbreak{}\hskip.1em \RB{17.179}{8.308}{0.502}{%
\RG{0.000}{-0.002}{10.9091}{mi}{66}\RG{8.814}{-0.002}{10.9091}{mi}{65}%
}\hskip.1em\allowbreak{} 。

\section*{2014A10正交矩阵与秩的稳定性}
题目设\allowbreak{}\hskip.1em \RB{6.729}{5.324}{0.618}{%
\RG{0.000}{0.002}{10.9091}{mi}{110}%
}\hskip.1em\allowbreak{} 阶实方阵\allowbreak{}\hskip.1em \RB{8.365}{8.313}{0.500}{%
\RG{0.000}{0.002}{10.9091}{mi}{65}%
}\hskip.1em\allowbreak{} 满足\allowbreak{}\hskip.1em \RB{44.287}{8.730}{2.210}{%
\RG{0.000}{0.002}{10.9091}{mi}{65}\RG{8.181}{0.002}{10.9091}{mi}{65}\RG{16.363}{3.962}{7.9701}{sy}{48}\RG{22.453}{0.002}{10.9091}{rm}{61}\RG{33.968}{0.003}{10.9091}{mi}{73}\RG{38.763}{-1.634}{7.9701}{mi}{110}%
}\hskip.1em\allowbreak{} （即\allowbreak{}\hskip.1em \RB{8.365}{8.313}{0.500}{%
\RG{0.000}{0.002}{10.9091}{mi}{65}%
}\hskip.1em\allowbreak{} 为\allowbreak{}\hskip.1em \RB{6.729}{5.324}{0.618}{%
\RG{0.000}{0.002}{10.9091}{mi}{110}%
}\hskip.1em\allowbreak{} 阶正交矩阵）。证明

\[
\RB{156.608}{10.077}{4.304}{%
\RG{98.389}{8.840}{10.9087}{ex}{0}\RG{149.014}{8.840}{10.9087}{ex}{1}\RG{144.120}{4.507}{7.9701}{rm}{50}\RT{0.000}{0.003}{10.9091}{r}\RT{4.276}{0.003}{10.9091}{a}\RT{9.731}{0.003}{10.9091}{n}\RT{15.796}{0.003}{10.9091}{k}\RG{21.555}{0.003}{10.9091}{rm}{40}\RG{25.798}{0.003}{10.9091}{mi}{73}\RG{38.955}{0.003}{10.9091}{sy}{0}\RG{49.865}{0.003}{10.9091}{mi}{65}\RG{58.046}{0.003}{10.9091}{rm}{41}\RG{65.300}{0.003}{10.9091}{rm}{61}\RT{76.832}{0.003}{10.9091}{r}\RT{81.108}{0.003}{10.9091}{a}\RT{86.563}{0.003}{10.9091}{n}\RT{92.628}{0.003}{10.9091}{k}\RG{103.388}{0.003}{10.9091}{rm}{40}\RG{107.630}{0.003}{10.9091}{mi}{73}\RG{120.787}{0.003}{10.9091}{sy}{0}\RG{131.696}{0.003}{10.9091}{mi}{65}\RG{139.878}{0.003}{10.9091}{rm}{41}\RG{154.014}{0.003}{10.9091}{mi}{58}\RG{30.593}{-1.634}{7.9701}{mi}{110}\RG{112.426}{-1.634}{7.9701}{mi}{110}%
}
\]
解：令\allowbreak{}\hskip.1em \RB{54.866}{8.307}{2.215}{%
\RG{0.000}{-0.003}{10.9091}{mi}{84}\RG{10.920}{-0.003}{10.9091}{rm}{61}\RG{22.435}{-0.003}{10.9091}{mi}{73}\RG{27.231}{-1.639}{7.9701}{mi}{110}\RG{35.592}{-0.003}{10.9091}{sy}{0}\RG{46.501}{-0.003}{10.9091}{mi}{65}%
}\hskip.1em\allowbreak{} 。显然ker\allowbreak{}\hskip.1em \RB{18.991}{7.882}{1.998}{%
\RG{0.000}{-0.003}{10.9091}{mi}{84}\RG{10.920}{-0.003}{10.9091}{sy}{18}%
}\hskip.1em\allowbreak{} ker\allowbreak{}\hskip.1em \RB{12.246}{9.526}{0.503}{%
\RG{0.000}{-0.003}{10.9091}{mi}{84}\RG{7.890}{3.956}{7.9701}{rm}{50}%
}\hskip.1em\allowbreak{} ，只需证明反向包含。

设\allowbreak{}\hskip.1em \RB{39.065}{9.527}{0.741}{%
\RG{0.000}{-0.001}{10.9091}{mi}{84}\RG{7.890}{3.957}{7.9701}{rm}{50}\RG{12.784}{-0.001}{10.9091}{mi}{120}\RG{22.049}{-0.001}{10.9091}{rm}{61}\RG{33.547}{-0.001}{10.9091}{rm}{48}%
}\hskip.1em\allowbreak{} ，令\allowbreak{}\hskip.1em \RB{34.409}{7.884}{2.737}{%
\RG{0.000}{-0.001}{10.9091}{mi}{121}\RG{8.770}{-0.001}{10.9091}{rm}{61}\RG{20.284}{-0.001}{10.9091}{mi}{84}\RG{28.160}{-0.001}{10.9091}{mi}{120}%
}\hskip.1em\allowbreak{} ，则\allowbreak{}\hskip.1em \RB{33.676}{7.884}{2.737}{%
\RG{0.000}{-0.001}{10.9091}{mi}{84}\RG{7.876}{-0.001}{10.9091}{mi}{121}\RG{16.660}{-0.001}{10.9091}{rm}{61}\RG{28.158}{-0.001}{10.9091}{rm}{48}%
}\hskip.1em\allowbreak{} ，即\allowbreak{}\hskip.1em \RB{34.311}{8.309}{2.737}{%
\RG{0.000}{-0.001}{10.9091}{mi}{65}\RG{8.181}{-0.001}{10.9091}{mi}{121}\RG{16.951}{-0.001}{10.9091}{rm}{61}\RG{28.466}{-0.001}{10.9091}{mi}{121}%
}\hskip.1em\allowbreak{} 。由\allowbreak{}\hskip.1em \RB{44.288}{8.726}{2.213}{%
\RG{0.000}{-0.001}{10.9091}{mi}{65}\RG{8.181}{-0.001}{10.9091}{mi}{65}\RG{16.363}{3.957}{7.9701}{sy}{48}\RG{22.454}{-0.001}{10.9091}{rm}{61}\RG{33.969}{-0.001}{10.9091}{mi}{73}\RG{38.764}{-1.637}{7.9701}{mi}{110}%
}\hskip.1em\allowbreak{} 可知\allowbreak{}\hskip.1em \RB{44.898}{9.527}{0.501}{%
\RG{0.000}{-0.001}{10.9091}{mi}{65}\RG{8.182}{3.957}{7.9701}{sy}{48}\RG{14.273}{-0.001}{10.9091}{rm}{61}\RG{25.787}{-0.001}{10.9091}{mi}{65}\RG{33.969}{3.957}{7.9701}{sy}{0}\RG{40.786}{3.957}{7.9701}{rm}{49}%
}\hskip.1em\allowbreak{} ，故\allowbreak{}\hskip.1em \RB{37.371}{8.726}{2.737}{%
\RG{0.000}{-0.001}{10.9091}{mi}{65}\RG{8.182}{3.957}{7.9701}{sy}{48}\RG{11.242}{-0.001}{10.9091}{mi}{121}\RG{20.012}{-0.001}{10.9091}{rm}{61}\RG{31.526}{-0.001}{10.9091}{mi}{121}%
}\hskip.1em\allowbreak{} ，从而\allowbreak{}\hskip.1em \RB{36.736}{8.726}{2.737}{%
\RG{0.000}{-0.001}{10.9091}{mi}{84}\RG{7.890}{3.957}{7.9701}{sy}{48}\RG{10.951}{-0.001}{10.9091}{mi}{121}\RG{19.720}{-0.001}{10.9091}{rm}{61}\RG{31.218}{-0.001}{10.9091}{rm}{48}%
}\hskip.1em\allowbreak{} 。因此

\[
\RB{123.618}{9.279}{3.222}{%
\RG{5.740}{4.510}{7.9701}{sy}{48}\RG{51.694}{4.510}{7.9701}{sy}{48}\RG{81.273}{4.510}{7.9701}{sy}{48}\RG{92.224}{4.510}{7.9701}{sy}{48}\RG{0.000}{0.005}{10.9091}{mi}{121}\RG{8.800}{0.005}{10.9091}{mi}{121}\RG{17.571}{0.005}{10.9091}{rm}{61}\RG{29.069}{0.005}{10.9091}{rm}{40}\RG{33.327}{0.005}{10.9091}{mi}{84}\RG{41.203}{0.005}{10.9091}{mi}{120}\RG{47.452}{0.005}{10.9091}{rm}{41}\RG{54.755}{0.005}{10.9091}{mi}{121}\RG{63.524}{0.005}{10.9091}{rm}{61}\RG{75.040}{0.005}{10.9091}{mi}{120}\RG{84.335}{0.005}{10.9091}{mi}{84}\RG{95.285}{0.005}{10.9091}{mi}{121}\RG{104.054}{0.005}{10.9091}{rm}{61}\RG{115.552}{0.005}{10.9091}{rm}{48}\RG{121.024}{0.005}{10.9091}{mi}{58}%
}
\]
因为\allowbreak{}\hskip.1em \RB{5.845}{5.322}{2.736}{%
\RG{0.000}{0.000}{10.9091}{mi}{121}%
}\hskip.1em\allowbreak{} 是实向量，\allowbreak{}\hskip.1em \RB{76.729}{9.530}{3.848}{%
\RG{0.000}{0.000}{10.9091}{mi}{121}\RG{5.740}{3.959}{7.9701}{sy}{48}\RG{8.800}{0.000}{10.9091}{mi}{121}\RG{17.571}{0.000}{10.9091}{rm}{61}\RG{29.085}{8.182}{10.9087}{ex}{80}\RG{40.600}{-3.272}{7.9701}{mi}{105}\RG{46.050}{0.000}{10.9091}{mi}{121}\RG{51.399}{-2.839}{7.9701}{mi}{105}\RG{51.789}{3.959}{7.9701}{rm}{50}\RG{59.714}{0.000}{10.9091}{rm}{61}\RG{71.211}{0.000}{10.9091}{rm}{48}%
}\hskip.1em\allowbreak{} 蕴含\allowbreak{}\hskip.1em \RB{25.785}{7.765}{2.736}{%
\RG{0.000}{0.000}{10.9091}{mi}{121}\RG{8.770}{0.000}{10.9091}{rm}{61}\RG{20.267}{0.000}{10.9091}{rm}{48}%
}\hskip.1em\allowbreak{} 。于是\allowbreak{}\hskip.1em \RB{34.170}{7.885}{0.740}{%
\RG{0.000}{0.000}{10.9091}{mi}{84}\RG{7.876}{0.000}{10.9091}{mi}{120}\RG{17.155}{0.000}{10.9091}{rm}{61}\RG{28.652}{0.000}{10.9091}{rm}{48}%
}\hskip.1em\allowbreak{} ，故ker\allowbreak{}\hskip.1em \RB{23.885}{9.530}{1.994}{%
\RG{0.000}{0.000}{10.9091}{mi}{84}\RG{7.890}{3.959}{7.9701}{rm}{50}\RG{15.815}{0.000}{10.9091}{sy}{18}%
}\hskip.1em\allowbreak{} ker\allowbreak{}\hskip.1em \RB{8.180}{7.885}{0.500}{%
\RG{0.000}{0.000}{10.9091}{mi}{84}%
}\hskip.1em\allowbreak{} 。

综上ker\allowbreak{}\hskip.1em \RB{24.180}{9.531}{0.500}{%
\RG{0.000}{0.002}{10.9091}{mi}{84}\RG{7.890}{3.960}{7.9701}{rm}{50}\RG{15.815}{0.002}{10.9091}{rm}{61}%
}\hskip.1em\allowbreak{} ker\allowbreak{}\hskip.1em \RB{8.180}{7.887}{0.500}{%
\RG{0.000}{0.002}{10.9091}{mi}{84}%
}\hskip.1em\allowbreak{} ，由秩–零度定理即得所需结论。\allowbreak{}\hskip.1em \RB{8.365}{7.996}{0.500}{%
\RG{0.000}{0.002}{10.9091}{am}{3}%
}\hskip.1em\allowbreak{}

小结：证明两矩阵同秩可转为证明同核。这里正交性给出ker\allowbreak{}\hskip.1em \RB{19.862}{7.886}{0.500}{%
\RG{0.000}{0.001}{10.9091}{mi}{84}\RG{11.497}{0.001}{10.9091}{rm}{61}%
}\hskip.1em\allowbreak{} ker\allowbreak{}\hskip.1em \RB{10.673}{8.729}{0.500}{%
\RG{0.000}{0.001}{10.9091}{mi}{84}\RG{7.890}{3.961}{7.9701}{sy}{48}%
}\hskip.1em\allowbreak{} ，再用实内积的正定性排除非零向量

\[
\RB{63.453}{9.297}{2.677}{%
\RG{0.000}{0.004}{10.9091}{mi}{84}\RG{7.876}{0.004}{10.9091}{mi}{120}\RG{17.155}{0.004}{10.9091}{sy}{50}\RT{27.457}{0.004}{10.9091}{k}\RT{33.217}{0.004}{10.9091}{e}\RT{38.060}{0.004}{10.9091}{r}\RG{44.154}{0.004}{10.9091}{mi}{84}\RT{52.045}{0.004}{10.9091}{。}%
}
\]
% Original page 4.
\section*{2014A11秩可加的幂等矩阵之和}
题目设\allowbreak{}\hskip.1em \RB{6.729}{5.317}{0.624}{%
\RG{0.000}{-0.004}{10.9091}{mi}{110}%
}\hskip.1em\allowbreak{} 阶方阵\allowbreak{}\hskip.1em \RB{21.761}{8.306}{2.610}{%
\RG{0.000}{-0.004}{10.9091}{mi}{65}\RG{8.181}{-0.004}{10.9091}{mi}{59}\RG{13.014}{-0.004}{10.9091}{mi}{66}%
}\hskip.1em\allowbreak{} 满足

\[
\RB{269.852}{10.071}{3.232}{%
\RG{38.821}{4.500}{7.9701}{rm}{50}\RG{0.000}{-0.004}{10.9091}{rm}{40}\RG{4.243}{-0.004}{10.9091}{mi}{65}\RG{14.849}{-0.004}{10.9091}{rm}{43}\RG{25.757}{-0.004}{10.9091}{mi}{66}\RG{34.579}{-0.004}{10.9091}{rm}{41}\RG{46.745}{-0.004}{10.9091}{rm}{61}\RG{58.260}{-0.004}{10.9091}{mi}{65}\RG{68.865}{-0.004}{10.9091}{rm}{43}\RG{79.773}{-0.004}{10.9091}{mi}{66}\RG{88.588}{-0.004}{10.9091}{mi}{59}\RT{115.261}{-0.004}{10.9091}{r}\RT{119.537}{-0.004}{10.9091}{a}\RT{124.992}{-0.004}{10.9091}{n}\RT{131.057}{-0.004}{10.9091}{k}\RG{136.816}{-0.004}{10.9091}{rm}{40}\RG{141.059}{-0.004}{10.9091}{mi}{65}\RG{151.665}{-0.004}{10.9091}{rm}{43}\RG{162.573}{-0.004}{10.9091}{mi}{66}\RG{171.395}{-0.004}{10.9091}{rm}{41}\RG{178.650}{-0.004}{10.9091}{rm}{61}\RT{190.183}{-0.004}{10.9091}{r}\RT{194.459}{-0.004}{10.9091}{a}\RT{199.913}{-0.004}{10.9091}{n}\RT{205.978}{-0.004}{10.9091}{k}\RG{213.556}{-0.004}{10.9091}{mi}{65}\RG{224.161}{-0.004}{10.9091}{rm}{43}\RT{235.071}{-0.004}{10.9091}{r}\RT{239.347}{-0.004}{10.9091}{a}\RT{244.801}{-0.004}{10.9091}{n}\RT{250.866}{-0.004}{10.9091}{k}\RG{258.444}{-0.004}{10.9091}{mi}{66}\RG{267.258}{-0.004}{10.9091}{mi}{58}%
}
\]
证明：\allowbreak{}\hskip.1em \RB{161.513}{9.525}{2.609}{%
\RG{0.000}{-0.003}{10.9091}{mi}{65}\RG{8.182}{3.955}{7.9701}{rm}{50}\RG{16.106}{-0.003}{10.9091}{rm}{61}\RG{27.620}{-0.003}{10.9091}{mi}{65}\RG{35.802}{-0.003}{10.9091}{mi}{59}\RG{44.267}{-0.003}{10.9091}{mi}{66}\RG{53.104}{3.955}{7.9701}{rm}{50}\RG{61.028}{-0.003}{10.9091}{rm}{61}\RG{72.542}{-0.003}{10.9091}{mi}{66}\RG{81.356}{-0.003}{10.9091}{mi}{59}\RG{89.822}{-0.003}{10.9091}{mi}{65}\RG{98.003}{-0.003}{10.9091}{mi}{66}\RG{109.878}{-0.003}{10.9091}{rm}{61}\RG{121.393}{-0.003}{10.9091}{mi}{66}\RG{130.207}{-0.003}{10.9091}{mi}{65}\RG{141.426}{-0.003}{10.9091}{rm}{61}\RG{152.941}{-0.003}{10.9091}{mi}{79}%
}\hskip.1em\allowbreak{} 。

解：令\allowbreak{}\hskip.1em \RB{53.383}{8.311}{1.405}{%
\RG{0.000}{0.000}{10.9091}{mi}{67}\RG{11.607}{0.000}{10.9091}{rm}{61}\RG{23.122}{0.000}{10.9091}{mi}{65}\RG{33.727}{0.000}{10.9091}{rm}{43}\RG{44.636}{0.000}{10.9091}{mi}{66}%
}\hskip.1em\allowbreak{} ，\allowbreak{}\hskip.1em \RB{20.033}{7.951}{0.740}{%
\RG{0.000}{0.000}{10.9091}{mi}{85}\RG{11.668}{0.000}{10.9091}{rm}{61}%
}\hskip.1em\allowbreak{} Im\allowbreak{}\hskip.1em \RB{8.365}{8.311}{0.500}{%
\RG{0.000}{0.000}{10.9091}{mi}{65}%
}\hskip.1em\allowbreak{} ，\allowbreak{}\hskip.1em \RB{23.213}{7.951}{0.740}{%
\RG{0.000}{0.000}{10.9091}{mi}{87}\RG{14.848}{0.000}{10.9091}{rm}{61}%
}\hskip.1em\allowbreak{} Im\allowbreak{}\hskip.1em \RB{8.747}{7.951}{0.500}{%
\RG{0.000}{0.000}{10.9091}{mi}{66}%
}\hskip.1em\allowbreak{} 。由Im\allowbreak{}\hskip.1em \RB{57.024}{8.191}{1.994}{%
\RG{0.000}{0.000}{10.9091}{mi}{67}\RG{11.607}{0.000}{10.9091}{sy}{18}\RG{23.121}{0.000}{10.9091}{mi}{85}\RG{34.184}{0.000}{10.9091}{rm}{43}\RG{45.092}{0.000}{10.9091}{mi}{87}%
}\hskip.1em\allowbreak{} ，有

\[
\RB{273.753}{9.292}{3.227}{%
\RT{0.000}{-0.000}{10.9091}{r}\RT{4.276}{-0.000}{10.9091}{a}\RT{9.731}{-0.000}{10.9091}{n}\RT{15.796}{-0.000}{10.9091}{k}\RG{23.373}{-0.000}{10.9091}{mi}{67}\RG{34.981}{-0.000}{10.9091}{sy}{20}\RT{46.495}{-0.000}{10.9091}{d}\RT{52.561}{-0.000}{10.9091}{i}\RT{55.593}{-0.000}{10.9091}{m}\RG{64.680}{-0.000}{10.9091}{rm}{40}\RG{68.922}{-0.000}{10.9091}{mi}{85}\RG{79.984}{-0.000}{10.9091}{rm}{43}\RG{90.893}{-0.000}{10.9091}{mi}{87}\RG{102.711}{-0.000}{10.9091}{rm}{41}\RG{109.965}{-0.000}{10.9091}{rm}{61}\RT{121.497}{-0.000}{10.9091}{r}\RT{125.773}{-0.000}{10.9091}{a}\RT{131.227}{-0.000}{10.9091}{n}\RT{137.293}{-0.000}{10.9091}{k}\RG{144.871}{-0.000}{10.9091}{mi}{65}\RG{155.477}{-0.000}{10.9091}{rm}{43}\RT{166.386}{-0.000}{10.9091}{r}\RT{170.662}{-0.000}{10.9091}{a}\RT{176.116}{-0.000}{10.9091}{n}\RT{182.182}{-0.000}{10.9091}{k}\RG{189.759}{-0.000}{10.9091}{mi}{66}\RG{201.005}{-0.000}{10.9091}{sy}{0}\RT{211.913}{-0.000}{10.9091}{d}\RT{217.978}{-0.000}{10.9091}{i}\RT{221.011}{-0.000}{10.9091}{m}\RG{230.098}{-0.000}{10.9091}{rm}{40}\RG{234.340}{-0.000}{10.9091}{mi}{85}\RG{245.403}{-0.000}{10.9091}{sy}{92}\RG{255.098}{-0.000}{10.9091}{mi}{87}\RG{266.916}{-0.000}{10.9091}{rm}{41}\RG{271.158}{-0.000}{10.9091}{mi}{58}%
}
\]
结合rank\allowbreak{}\hskip.1em \RB{19.973}{8.192}{0.739}{%
\RG{0.000}{0.001}{10.9091}{mi}{67}\RG{11.608}{0.001}{10.9091}{rm}{61}%
}\hskip.1em\allowbreak{} rank\allowbreak{}\hskip.1em \RB{18.971}{8.312}{1.404}{%
\RG{0.000}{0.001}{10.9091}{mi}{65}\RG{10.606}{0.001}{10.9091}{rm}{43}%
}\hskip.1em\allowbreak{} rank\allowbreak{}\hskip.1em \RB{8.747}{7.952}{0.500}{%
\RG{0.000}{0.001}{10.9091}{mi}{66}%
}\hskip.1em\allowbreak{} ，得到\allowbreak{}\hskip.1em \RB{63.187}{8.683}{3.226}{%
\RG{0.000}{0.001}{10.9091}{mi}{85}\RG{11.061}{0.001}{10.9091}{sy}{92}\RG{20.758}{0.001}{10.9091}{mi}{87}\RG{35.606}{0.001}{10.9091}{rm}{61}\RG{47.121}{0.001}{10.9091}{sy}{102}\RG{52.575}{0.001}{10.9091}{rm}{48}\RG{58.029}{0.001}{10.9091}{sy}{103}%
}\hskip.1em\allowbreak{} ，且Im\allowbreak{}\hskip.1em \RB{57.025}{8.192}{1.404}{%
\RG{0.000}{0.001}{10.9091}{mi}{67}\RG{11.608}{0.001}{10.9091}{rm}{61}\RG{23.122}{0.001}{10.9091}{mi}{85}\RG{34.184}{0.001}{10.9091}{sy}{8}\RG{45.093}{0.001}{10.9091}{mi}{87}%
}\hskip.1em\allowbreak{} 。

由于\allowbreak{}\hskip.1em \RB{36.806}{9.532}{0.738}{%
\RG{0.000}{0.002}{10.9091}{mi}{67}\RG{8.577}{3.962}{7.9701}{rm}{50}\RG{16.501}{0.002}{10.9091}{rm}{61}\RG{28.016}{0.002}{10.9091}{mi}{67}%
}\hskip.1em\allowbreak{} ，\allowbreak{}\hskip.1em \RB{8.791}{8.193}{0.738}{%
\RG{0.000}{0.002}{10.9091}{mi}{67}%
}\hskip.1em\allowbreak{} 在Im\allowbreak{}\hskip.1em \RB{8.791}{8.193}{0.738}{%
\RG{0.000}{0.002}{10.9091}{mi}{67}%
}\hskip.1em\allowbreak{} 上为恒等变换：若\allowbreak{}\hskip.1em \RB{35.097}{8.193}{2.734}{%
\RG{0.000}{0.002}{10.9091}{mi}{121}\RG{8.771}{0.002}{10.9091}{rm}{61}\RG{20.285}{0.002}{10.9091}{mi}{67}\RG{28.849}{0.002}{10.9091}{mi}{120}%
}\hskip.1em\allowbreak{} ，则\allowbreak{}\hskip.1em \RB{98.312}{9.532}{2.734}{%
\RG{0.000}{0.002}{10.9091}{mi}{67}\RG{8.563}{0.002}{10.9091}{mi}{121}\RG{17.346}{0.002}{10.9091}{rm}{61}\RG{28.861}{0.002}{10.9091}{mi}{67}\RG{37.438}{3.962}{7.9701}{rm}{50}\RG{42.332}{0.002}{10.9091}{mi}{120}\RG{51.597}{0.002}{10.9091}{rm}{61}\RG{63.111}{0.002}{10.9091}{mi}{67}\RG{71.675}{0.002}{10.9091}{mi}{120}\RG{80.953}{0.002}{10.9091}{rm}{61}\RG{92.467}{0.002}{10.9091}{mi}{121}%
}\hskip.1em\allowbreak{} 。

对任意\allowbreak{}\hskip.1em \RB{38.812}{7.955}{2.732}{%
\RG{0.000}{0.004}{10.9091}{mi}{121}\RG{8.771}{0.004}{10.9091}{sy}{50}\RG{19.073}{0.004}{10.9091}{mi}{85}\RG{30.741}{0.004}{10.9091}{sy}{18}%
}\hskip.1em\allowbreak{} Im\allowbreak{}\hskip.1em \RB{8.791}{8.195}{0.736}{%
\RG{0.000}{0.004}{10.9091}{mi}{67}%
}\hskip.1em\allowbreak{} ，由\allowbreak{}\hskip.1em \RB{34.707}{8.195}{2.732}{%
\RG{0.000}{0.004}{10.9091}{mi}{67}\RG{8.563}{0.004}{10.9091}{mi}{121}\RG{17.347}{0.004}{10.9091}{rm}{61}\RG{28.862}{0.004}{10.9091}{mi}{121}%
}\hskip.1em\allowbreak{} 得

\[
\RB{91.966}{8.678}{3.230}{%
\RG{0.000}{-0.003}{10.9091}{rm}{40}\RG{4.242}{-0.003}{10.9091}{mi}{65}\RG{12.423}{-0.003}{10.9091}{mi}{121}\RG{20.587}{-0.003}{10.9091}{sy}{0}\RG{31.495}{-0.003}{10.9091}{mi}{121}\RG{37.235}{-0.003}{10.9091}{rm}{41}\RG{43.890}{-0.003}{10.9091}{rm}{43}\RG{54.810}{-0.003}{10.9091}{mi}{66}\RG{63.625}{-0.003}{10.9091}{mi}{121}\RG{72.402}{-0.003}{10.9091}{rm}{61}\RG{83.900}{-0.003}{10.9091}{rm}{48}\RG{89.372}{-0.003}{10.9091}{mi}{58}%
}
\]
其中\allowbreak{}\hskip.1em \RB{54.160}{8.313}{2.734}{%
\RG{0.000}{0.002}{10.9091}{mi}{65}\RG{8.181}{0.002}{10.9091}{mi}{121}\RG{15.867}{0.002}{10.9091}{sy}{0}\RG{26.298}{0.002}{10.9091}{mi}{121}\RG{35.067}{0.002}{10.9091}{sy}{50}\RG{45.370}{0.002}{10.9091}{mi}{85}%
}\hskip.1em\allowbreak{} ，\allowbreak{}\hskip.1em \RB{39.826}{7.953}{2.734}{%
\RG{0.000}{0.002}{10.9091}{mi}{66}\RG{8.814}{0.002}{10.9091}{mi}{121}\RG{17.592}{0.002}{10.9091}{sy}{50}\RG{27.894}{0.002}{10.9091}{mi}{87}%
}\hskip.1em\allowbreak{} ；因\allowbreak{}\hskip.1em \RB{62.231}{8.684}{3.225}{%
\RG{0.000}{0.002}{10.9091}{mi}{85}\RG{10.583}{0.002}{10.9091}{sy}{92}\RG{19.801}{0.002}{10.9091}{mi}{87}\RG{34.649}{0.002}{10.9091}{rm}{61}\RG{46.165}{0.002}{10.9091}{sy}{102}\RG{51.619}{0.002}{10.9091}{rm}{48}\RG{57.073}{0.002}{10.9091}{sy}{103}%
}\hskip.1em\allowbreak{} ，必有\allowbreak{}\hskip.1em \RB{34.312}{8.313}{2.734}{%
\RG{0.000}{0.002}{10.9091}{mi}{65}\RG{8.182}{0.002}{10.9091}{mi}{121}\RG{16.952}{0.002}{10.9091}{rm}{61}\RG{28.467}{0.002}{10.9091}{mi}{121}%
}\hskip.1em\allowbreak{} 、\allowbreak{}\hskip.1em \RB{34.607}{7.953}{2.734}{%
\RG{0.000}{0.002}{10.9091}{mi}{66}\RG{8.814}{0.002}{10.9091}{mi}{121}\RG{17.592}{0.002}{10.9091}{rm}{61}\RG{29.089}{0.002}{10.9091}{rm}{48}%
}\hskip.1em\allowbreak{} 。取\allowbreak{}\hskip.1em \RB{34.716}{8.313}{2.734}{%
\RG{0.000}{0.002}{10.9091}{mi}{121}\RG{8.770}{0.002}{10.9091}{rm}{61}\RG{20.285}{0.002}{10.9091}{mi}{65}\RG{28.467}{0.002}{10.9091}{mi}{120}%
}\hskip.1em\allowbreak{} ，即得\allowbreak{}\hskip.1em \RB{48.285}{9.532}{0.618}{%
\RG{0.000}{0.002}{10.9091}{mi}{65}\RG{8.182}{3.962}{7.9701}{rm}{50}\RG{13.076}{0.002}{10.9091}{mi}{120}\RG{22.340}{0.002}{10.9091}{rm}{61}\RG{33.855}{0.002}{10.9091}{mi}{65}\RG{42.037}{0.002}{10.9091}{mi}{120}%
}\hskip.1em\allowbreak{} 、\allowbreak{}\hskip.1em \RB{43.284}{8.313}{0.738}{%
\RG{0.000}{0.002}{10.9091}{mi}{66}\RG{8.814}{0.002}{10.9091}{mi}{65}\RG{16.996}{0.002}{10.9091}{mi}{120}\RG{26.268}{0.002}{10.9091}{rm}{61}\RG{37.766}{0.002}{10.9091}{rm}{48}%
}\hskip.1em\allowbreak{} ，所以\allowbreak{}\hskip.1em \RB{35.986}{9.524}{0.505}{%
\RG{0.000}{-0.005}{10.9091}{mi}{65}\RG{8.182}{3.954}{7.9701}{rm}{50}\RG{16.105}{-0.005}{10.9091}{rm}{61}\RG{27.620}{-0.005}{10.9091}{mi}{65}%
}\hskip.1em\allowbreak{} 、\allowbreak{}\hskip.1em \RB{40.120}{8.305}{0.745}{%
\RG{0.000}{-0.005}{10.9091}{mi}{66}\RG{8.814}{-0.005}{10.9091}{mi}{65}\RG{20.033}{-0.005}{10.9091}{rm}{61}\RG{31.548}{-0.005}{10.9091}{mi}{79}%
}\hskip.1em\allowbreak{} 。

同理，对\allowbreak{}\hskip.1em \RB{30.818}{7.947}{0.940}{%
\RG{0.000}{-0.003}{10.9091}{mi}{122}\RG{8.583}{-0.003}{10.9091}{sy}{50}\RG{18.886}{-0.003}{10.9091}{mi}{87}%
}\hskip.1em\allowbreak{} ，等式\allowbreak{}\hskip.1em \RB{88.851}{8.678}{3.230}{%
\RG{0.000}{-0.003}{10.9091}{mi}{65}\RG{8.181}{-0.003}{10.9091}{mi}{122}\RG{16.158}{-0.003}{10.9091}{rm}{43}\RG{27.056}{-0.003}{10.9091}{rm}{40}\RG{31.309}{-0.003}{10.9091}{mi}{66}\RG{40.123}{-0.003}{10.9091}{mi}{122}\RG{48.108}{-0.003}{10.9091}{sy}{0}\RG{59.017}{-0.003}{10.9091}{mi}{122}\RG{64.570}{-0.003}{10.9091}{rm}{41}\RG{71.824}{-0.003}{10.9091}{rm}{61}\RG{83.333}{-0.003}{10.9091}{rm}{48}%
}\hskip.1em\allowbreak{} 给出\allowbreak{}\hskip.1em \RB{33.781}{8.307}{0.743}{%
\RG{0.000}{-0.003}{10.9091}{mi}{65}\RG{8.181}{-0.003}{10.9091}{mi}{122}\RG{16.765}{-0.003}{10.9091}{rm}{61}\RG{28.263}{-0.003}{10.9091}{rm}{48}%
}\hskip.1em\allowbreak{} 、\allowbreak{}\hskip.1em \RB{34.513}{7.947}{0.623}{%
\RG{0.000}{-0.003}{10.9091}{mi}{66}\RG{8.814}{-0.003}{10.9091}{mi}{122}\RG{17.404}{-0.003}{10.9091}{rm}{61}\RG{28.919}{-0.003}{10.9091}{mi}{122}%
}\hskip.1em\allowbreak{} 。取\allowbreak{}\hskip.1em \RB{35.160}{7.947}{0.623}{%
\RG{0.000}{-0.003}{10.9091}{mi}{122}\RG{8.582}{-0.003}{10.9091}{rm}{61}\RG{20.097}{-0.003}{10.9091}{mi}{66}\RG{28.911}{-0.003}{10.9091}{mi}{120}%
}\hskip.1em\allowbreak{} ，得到\allowbreak{}\hskip.1em \RB{40.120}{8.307}{0.743}{%
\RG{0.000}{-0.003}{10.9091}{mi}{65}\RG{8.181}{-0.003}{10.9091}{mi}{66}\RG{20.033}{-0.003}{10.9091}{rm}{61}\RG{31.548}{-0.003}{10.9091}{mi}{79}%
}\hskip.1em\allowbreak{} 、\allowbreak{}\hskip.1em \RB{37.007}{9.525}{0.503}{%
\RG{0.000}{-0.003}{10.9091}{mi}{66}\RG{8.822}{3.955}{7.9701}{rm}{50}\RG{16.745}{-0.003}{10.9091}{rm}{61}\RG{28.260}{-0.003}{10.9091}{mi}{66}%
}\hskip.1em\allowbreak{} 。\allowbreak{}\hskip.1em \RB{8.365}{7.991}{0.503}{%
\RG{0.000}{-0.003}{10.9091}{am}{3}%
}\hskip.1em\allowbreak{}

小结：秩的加法取等号，迫使两个像空间构成直和。幂等变换在自己的像上为恒等变换，按直和分量比较即可得到四个结论。

\section*{2014A12零乘积条件下的秩可加性}
题目设\allowbreak{}\hskip.1em \RB{6.729}{5.318}{0.623}{%
\RG{0.000}{-0.003}{10.9091}{mi}{110}%
}\hskip.1em\allowbreak{} 阶方阵\allowbreak{}\hskip.1em \RB{21.761}{8.307}{2.609}{%
\RG{0.000}{-0.003}{10.9091}{mi}{65}\RG{8.181}{-0.003}{10.9091}{mi}{59}\RG{13.014}{-0.003}{10.9091}{mi}{66}%
}\hskip.1em\allowbreak{} 满足\allowbreak{}\hskip.1em \RB{71.668}{8.307}{0.743}{%
\RG{0.000}{-0.003}{10.9091}{mi}{65}\RG{8.181}{-0.003}{10.9091}{mi}{66}\RG{20.033}{-0.003}{10.9091}{rm}{61}\RG{31.548}{-0.003}{10.9091}{mi}{66}\RG{40.362}{-0.003}{10.9091}{mi}{65}\RG{51.581}{-0.003}{10.9091}{rm}{61}\RG{63.096}{-0.003}{10.9091}{mi}{79}%
}\hskip.1em\allowbreak{} 、rank\allowbreak{}\hskip.1em \RB{19.576}{8.307}{0.503}{%
\RG{0.000}{-0.003}{10.9091}{mi}{65}\RG{11.211}{-0.003}{10.9091}{rm}{61}%
}\hskip.1em\allowbreak{} rank\allowbreak{}\hskip.1em \RB{12.538}{9.525}{0.503}{%
\RG{0.000}{-0.003}{10.9091}{mi}{65}\RG{8.182}{3.955}{7.9701}{rm}{50}%
}\hskip.1em\allowbreak{} 。证明

\[
\RB{154.591}{9.289}{3.230}{%
\RT{0.000}{-0.003}{10.9091}{r}\RT{4.276}{-0.003}{10.9091}{a}\RT{9.731}{-0.003}{10.9091}{n}\RT{15.796}{-0.003}{10.9091}{k}\RG{21.555}{-0.003}{10.9091}{rm}{40}\RG{25.797}{-0.003}{10.9091}{mi}{65}\RG{36.403}{-0.003}{10.9091}{rm}{43}\RG{47.312}{-0.003}{10.9091}{mi}{66}\RG{56.133}{-0.003}{10.9091}{rm}{41}\RG{63.387}{-0.003}{10.9091}{rm}{61}\RT{74.920}{-0.003}{10.9091}{r}\RT{79.197}{-0.003}{10.9091}{a}\RT{84.651}{-0.003}{10.9091}{n}\RT{90.716}{-0.003}{10.9091}{k}\RG{98.295}{-0.003}{10.9091}{mi}{65}\RG{108.900}{-0.003}{10.9091}{rm}{43}\RT{119.808}{-0.003}{10.9091}{r}\RT{124.085}{-0.003}{10.9091}{a}\RT{129.539}{-0.003}{10.9091}{n}\RT{135.604}{-0.003}{10.9091}{k}\RG{143.183}{-0.003}{10.9091}{mi}{66}\RG{151.997}{-0.003}{10.9091}{mi}{58}%
}
\]
解：由Im\allowbreak{}\hskip.1em \RB{24.176}{9.526}{1.998}{%
\RG{0.000}{-0.003}{10.9091}{mi}{65}\RG{8.182}{3.956}{7.9701}{rm}{50}\RG{16.106}{-0.003}{10.9091}{sy}{18}%
}\hskip.1em\allowbreak{} Im\allowbreak{}\hskip.1em \RB{8.365}{8.307}{0.503}{%
\RG{0.000}{-0.003}{10.9091}{mi}{65}%
}\hskip.1em\allowbreak{} 、ker\allowbreak{}\hskip.1em \RB{19.283}{8.307}{1.998}{%
\RG{0.000}{-0.003}{10.9091}{mi}{65}\RG{11.212}{-0.003}{10.9091}{sy}{18}%
}\hskip.1em\allowbreak{} ker\allowbreak{}\hskip.1em \RB{12.538}{9.526}{0.503}{%
\RG{0.000}{-0.003}{10.9091}{mi}{65}\RG{8.182}{3.956}{7.9701}{rm}{50}%
}\hskip.1em\allowbreak{} 及秩相等，得到

\[
\RB{163.944}{10.072}{2.618}{%
\RG{23.024}{4.501}{7.9701}{rm}{50}\RG{117.032}{4.501}{7.9701}{rm}{50}\RT{0.000}{-0.002}{10.9091}{I}\RT{3.938}{-0.002}{10.9091}{m}\RG{14.842}{-0.002}{10.9091}{mi}{65}\RG{30.948}{-0.002}{10.9091}{rm}{61}\RT{42.463}{-0.002}{10.9091}{I}\RT{46.401}{-0.002}{10.9091}{m}\RG{57.306}{-0.002}{10.9091}{mi}{65}\RG{65.487}{-0.002}{10.9091}{mi}{59}\RT{92.153}{-0.002}{10.9091}{k}\RT{97.913}{-0.002}{10.9091}{e}\RT{102.756}{-0.002}{10.9091}{r}\RG{108.850}{-0.002}{10.9091}{mi}{65}\RG{124.956}{-0.002}{10.9091}{rm}{61}\RT{136.470}{-0.002}{10.9091}{k}\RT{142.230}{-0.002}{10.9091}{e}\RT{147.074}{-0.002}{10.9091}{r}\RG{153.169}{-0.002}{10.9091}{mi}{65}\RG{161.350}{-0.002}{10.9091}{mi}{58}%
}
\]
若\allowbreak{}\hskip.1em \RB{15.629}{6.388}{2.738}{%
\RG{0.000}{-0.002}{10.9091}{mi}{121}\RG{8.770}{-0.002}{10.9091}{sy}{50}%
}\hskip.1em\allowbreak{} Im\allowbreak{}\hskip.1em \RB{15.692}{8.308}{0.742}{%
\RG{0.000}{-0.002}{10.9091}{mi}{65}\RG{8.526}{-0.002}{10.9091}{sy}{92}%
}\hskip.1em\allowbreak{} ker\allowbreak{}\hskip.1em \RB{8.365}{8.308}{0.502}{%
\RG{0.000}{-0.002}{10.9091}{mi}{65}%
}\hskip.1em\allowbreak{} ，则可写成\allowbreak{}\hskip.1em \RB{34.715}{8.308}{2.738}{%
\RG{0.000}{-0.002}{10.9091}{mi}{121}\RG{8.770}{-0.002}{10.9091}{rm}{61}\RG{20.284}{-0.002}{10.9091}{mi}{65}\RG{28.466}{-0.002}{10.9091}{mi}{120}%
}\hskip.1em\allowbreak{} ，且\allowbreak{}\hskip.1em \RB{67.790}{9.526}{2.738}{%
\RG{0.000}{-0.002}{10.9091}{rm}{48}\RG{8.476}{-0.002}{10.9091}{rm}{61}\RG{19.999}{-0.002}{10.9091}{mi}{65}\RG{28.181}{-0.002}{10.9091}{mi}{121}\RG{36.951}{-0.002}{10.9091}{rm}{61}\RG{48.465}{-0.002}{10.9091}{mi}{65}\RG{56.647}{3.956}{7.9701}{rm}{50}\RG{61.541}{-0.002}{10.9091}{mi}{120}%
}\hskip.1em\allowbreak{} ；由同核可知\allowbreak{}\hskip.1em \RB{34.462}{8.308}{0.742}{%
\RG{0.000}{-0.002}{10.9091}{mi}{65}\RG{8.181}{-0.002}{10.9091}{mi}{120}\RG{17.446}{-0.002}{10.9091}{rm}{61}\RG{28.944}{-0.002}{10.9091}{rm}{48}%
}\hskip.1em\allowbreak{} ，故\allowbreak{}\hskip.1em \RB{25.785}{7.763}{2.738}{%
\RG{0.000}{-0.002}{10.9091}{mi}{121}\RG{8.770}{-0.002}{10.9091}{rm}{61}\RG{20.267}{-0.002}{10.9091}{rm}{48}%
}\hskip.1em\allowbreak{} 。因此Im\allowbreak{}\hskip.1em \RB{15.691}{8.308}{0.742}{%
\RG{0.000}{-0.002}{10.9091}{mi}{65}\RG{8.525}{-0.002}{10.9091}{sy}{92}%
}\hskip.1em\allowbreak{} ker\allowbreak{}\hskip.1em \RB{38.793}{8.679}{3.229}{%
\RG{0.000}{-0.002}{10.9091}{mi}{65}\RG{11.211}{-0.002}{10.9091}{rm}{61}\RG{22.726}{-0.002}{10.9091}{sy}{102}\RG{28.180}{-0.002}{10.9091}{rm}{48}\RG{33.635}{-0.002}{10.9091}{sy}{103}%
}\hskip.1em\allowbreak{} 。

由\allowbreak{}\hskip.1em \RB{40.120}{8.311}{0.740}{%
\RG{0.000}{-0.000}{10.9091}{mi}{65}\RG{8.181}{-0.000}{10.9091}{mi}{66}\RG{20.033}{-0.000}{10.9091}{rm}{61}\RG{31.548}{-0.000}{10.9091}{mi}{79}%
}\hskip.1em\allowbreak{} ，有Im\allowbreak{}\hskip.1em \RB{19.922}{7.951}{1.995}{%
\RG{0.000}{-0.000}{10.9091}{mi}{66}\RG{11.852}{-0.000}{10.9091}{sy}{18}%
}\hskip.1em\allowbreak{} ker\allowbreak{}\hskip.1em \RB{8.365}{8.311}{0.500}{%
\RG{0.000}{-0.000}{10.9091}{mi}{65}%
}\hskip.1em\allowbreak{} ，于是Im\allowbreak{}\hskip.1em \RB{17.771}{8.311}{0.740}{%
\RG{0.000}{-0.000}{10.9091}{mi}{65}\RG{10.606}{-0.000}{10.9091}{sy}{92}%
}\hskip.1em\allowbreak{} Im\allowbreak{}\hskip.1em \RB{39.432}{8.681}{3.227}{%
\RG{0.000}{-0.000}{10.9091}{mi}{66}\RG{11.852}{-0.000}{10.9091}{rm}{61}\RG{23.366}{-0.000}{10.9091}{sy}{102}\RG{28.820}{-0.000}{10.9091}{rm}{48}\RG{34.274}{-0.000}{10.9091}{sy}{103}%
}\hskip.1em\allowbreak{} 。

再令\allowbreak{}\hskip.1em \RB{53.383}{8.312}{1.404}{%
\RG{0.000}{0.001}{10.9091}{mi}{67}\RG{11.607}{0.001}{10.9091}{rm}{61}\RG{23.121}{0.001}{10.9091}{mi}{65}\RG{33.727}{0.001}{10.9091}{rm}{43}\RG{44.636}{0.001}{10.9091}{mi}{66}%
}\hskip.1em\allowbreak{} 。由于\allowbreak{}\hskip.1em \RB{40.120}{8.312}{0.739}{%
\RG{0.000}{0.001}{10.9091}{mi}{66}\RG{8.814}{0.001}{10.9091}{mi}{65}\RG{20.033}{0.001}{10.9091}{rm}{61}\RG{31.548}{0.001}{10.9091}{mi}{79}%
}\hskip.1em\allowbreak{} ，有\allowbreak{}\hskip.1em \RB{43.840}{9.531}{0.739}{%
\RG{0.000}{0.001}{10.9091}{mi}{67}\RG{8.563}{0.001}{10.9091}{mi}{65}\RG{19.788}{0.001}{10.9091}{rm}{61}\RG{31.303}{0.001}{10.9091}{mi}{65}\RG{39.484}{3.960}{7.9701}{rm}{50}%
}\hskip.1em\allowbreak{} ，故

\[
\RB{158.246}{10.076}{3.225}{%
\RG{60.594}{4.506}{7.9701}{rm}{50}\RT{0.000}{0.002}{10.9091}{I}\RT{3.938}{0.002}{10.9091}{m}\RG{14.843}{0.002}{10.9091}{mi}{65}\RG{26.055}{0.002}{10.9091}{rm}{61}\RT{37.569}{0.002}{10.9091}{I}\RT{41.507}{0.002}{10.9091}{m}\RG{52.412}{0.002}{10.9091}{mi}{65}\RG{68.518}{0.002}{10.9091}{rm}{61}\RT{80.033}{0.002}{10.9091}{I}\RT{83.971}{0.002}{10.9091}{m}\RG{93.057}{0.002}{10.9091}{rm}{40}\RG{97.300}{0.002}{10.9091}{mi}{67}\RG{105.864}{0.002}{10.9091}{mi}{65}\RG{114.059}{0.002}{10.9091}{rm}{41}\RG{121.330}{0.002}{10.9091}{sy}{18}\RT{132.846}{0.002}{10.9091}{I}\RT{136.784}{0.002}{10.9091}{m}\RG{147.688}{0.002}{10.9091}{mi}{67}\RG{155.652}{0.002}{10.9091}{mi}{58}%
}
\]
对任意\allowbreak{}\hskip.1em \RB{6.249}{5.324}{0.618}{%
\RG{0.000}{0.002}{10.9091}{mi}{120}%
}\hskip.1em\allowbreak{} ，\allowbreak{}\hskip.1em \RB{82.052}{8.313}{0.934}{%
\RG{0.000}{0.002}{10.9091}{mi}{66}\RG{8.814}{0.002}{10.9091}{mi}{120}\RG{18.087}{0.002}{10.9091}{rm}{61}\RG{29.601}{0.002}{10.9091}{mi}{67}\RG{38.164}{0.002}{10.9091}{mi}{120}\RG{46.837}{0.002}{10.9091}{sy}{0}\RG{57.745}{0.002}{10.9091}{mi}{65}\RG{65.926}{0.002}{10.9091}{mi}{120}\RG{75.192}{0.002}{10.9091}{sy}{50}%
}\hskip.1em\allowbreak{} Im\allowbreak{}\hskip.1em \RB{8.791}{8.193}{0.738}{%
\RG{0.000}{0.002}{10.9091}{mi}{67}%
}\hskip.1em\allowbreak{} ，所以Im\allowbreak{}\hskip.1em \RB{19.922}{7.953}{1.992}{%
\RG{0.000}{0.002}{10.9091}{mi}{66}\RG{11.852}{0.002}{10.9091}{sy}{18}%
}\hskip.1em\allowbreak{} Im\allowbreak{}\hskip.1em \RB{8.791}{8.193}{0.738}{%
\RG{0.000}{0.002}{10.9091}{mi}{67}%
}\hskip.1em\allowbreak{} 。结合显然的反向包含，得到

Im\allowbreak{}\hskip.1em \RB{50.197}{8.684}{3.225}{%
\RG{0.000}{0.002}{10.9091}{rm}{40}\RG{4.242}{0.003}{10.9091}{mi}{65}\RG{14.848}{0.002}{10.9091}{rm}{43}\RG{25.756}{0.003}{10.9091}{mi}{66}\RG{34.578}{0.002}{10.9091}{rm}{41}\RG{41.832}{0.002}{10.9091}{rm}{61}%
}\hskip.1em\allowbreak{} Im\allowbreak{}\hskip.1em \RB{18.981}{8.313}{1.402}{%
\RG{0.000}{0.002}{10.9091}{mi}{65}\RG{10.605}{0.003}{10.9091}{sy}{8}%
}\hskip.1em\allowbreak{} Im\allowbreak{}\hskip.1em \RB{11.529}{7.954}{2.602}{%
\RG{0.000}{0.003}{10.9091}{mi}{66}\RG{8.814}{0.003}{10.9091}{mi}{59}%
}\hskip.1em\allowbreak{} 从而rank\allowbreak{}\hskip.1em \RB{50.198}{8.684}{3.225}{%
\RG{0.000}{0.002}{10.9091}{rm}{40}\RG{4.243}{0.003}{10.9091}{mi}{65}\RG{14.848}{0.002}{10.9091}{rm}{43}\RG{25.757}{0.003}{10.9091}{mi}{66}\RG{34.579}{0.002}{10.9091}{rm}{41}\RG{41.833}{0.002}{10.9091}{rm}{61}%
}\hskip.1em\allowbreak{} rank\allowbreak{}\hskip.1em \RB{18.971}{8.313}{1.402}{%
\RG{0.000}{0.002}{10.9091}{mi}{65}\RG{10.606}{0.003}{10.9091}{rm}{43}%
}\hskip.1em\allowbreak{} rank\allowbreak{}\hskip.1em \RB{15.125}{7.954}{0.500}{%
\RG{0.000}{0.003}{10.9091}{mi}{66}\RG{12.531}{0.003}{10.9091}{mi}{58}%
}\hskip.1em\allowbreak{}

小结：先证明像空间交为零，再证明和矩阵的像恰好等于两个像之和。条件rank\allowbreak{}\hskip.1em \RB{21.736}{8.309}{0.501}{%
\RG{0.000}{-0.001}{10.9091}{mi}{65}\RG{13.371}{-0.001}{10.9091}{rm}{61}%
}\hskip.1em\allowbreak{} rank\allowbreak{}\hskip.1em \RB{12.538}{9.527}{0.501}{%
\RG{0.000}{-0.001}{10.9091}{mi}{65}\RG{8.182}{3.957}{7.9701}{rm}{50}%
}\hskip.1em\allowbreak{} 正好保证Im\allowbreak{}\hskip.1em \RB{17.771}{8.313}{0.738}{%
\RG{0.000}{0.002}{10.9091}{mi}{65}\RG{10.606}{0.002}{10.9091}{sy}{92}%
}\hskip.1em\allowbreak{} ker\allowbreak{}\hskip.1em \RB{38.793}{8.684}{3.225}{%
\RG{0.000}{0.002}{10.9091}{mi}{65}\RG{11.211}{0.002}{10.9091}{rm}{61}\RG{22.726}{0.002}{10.9091}{sy}{102}\RG{28.180}{0.002}{10.9091}{rm}{48}\RG{33.635}{0.002}{10.9091}{sy}{103}%
}\hskip.1em\allowbreak{} ，并使Im\allowbreak{}\hskip.1em \RB{8.365}{8.313}{0.500}{%
\RG{0.000}{0.002}{10.9091}{mi}{65}%
}\hskip.1em\allowbreak{} 能由\allowbreak{}\hskip.1em \RB{47.186}{8.684}{3.225}{%
\RG{0.000}{0.002}{10.9091}{rm}{40}\RG{4.242}{0.002}{10.9091}{mi}{65}\RG{14.848}{0.002}{10.9091}{rm}{43}\RG{25.757}{0.002}{10.9091}{mi}{66}\RG{34.578}{0.002}{10.9091}{rm}{41}\RG{38.821}{0.002}{10.9091}{mi}{65}%
}\hskip.1em\allowbreak{} 生成。

% Original page 5.
\section*{2014A13幂零变换与循环向量}
题目设\allowbreak{}\hskip.1em \RB{8.889}{7.956}{0.735}{%
\RG{0.000}{0.005}{10.9091}{mi}{86}%
}\hskip.1em\allowbreak{} 是数域\allowbreak{}\hskip.1em \RB{8.889}{7.978}{0.500}{%
\RG{0.000}{0.005}{10.9091}{bb}{75}%
}\hskip.1em\allowbreak{} 上的\allowbreak{}\hskip.1em \RB{6.729}{5.327}{0.615}{%
\RG{0.000}{0.005}{10.9091}{mi}{110}%
}\hskip.1em\allowbreak{} 维线性空间，\allowbreak{}\hskip.1em \RB{7.252}{5.327}{2.873}{%
\RG{0.000}{0.005}{10.9091}{mi}{39}%
}\hskip.1em\allowbreak{} 是\allowbreak{}\hskip.1em \RB{8.889}{7.956}{0.735}{%
\RG{0.000}{0.005}{10.9091}{mi}{86}%
}\hskip.1em\allowbreak{} 上的幂零线性变换，满足rank\allowbreak{}\hskip.1em \RB{45.849}{7.771}{2.873}{%
\RG{0.000}{0.005}{10.9091}{mi}{39}\RG{10.166}{0.005}{10.9091}{rm}{61}\RG{21.680}{0.005}{10.9091}{mi}{110}\RG{30.261}{0.005}{10.9091}{sy}{0}\RG{40.779}{0.005}{10.9091}{rm}{49}%
}\hskip.1em\allowbreak{} 。求证：\allowbreak{}\hskip.1em \RB{8.889}{7.956}{0.735}{%
\RG{0.000}{0.005}{10.9091}{mi}{86}%
}\hskip.1em\allowbreak{} 是关于线性变换\allowbreak{}\hskip.1em \RB{7.252}{5.319}{2.880}{%
\RG{0.000}{-0.002}{10.9091}{mi}{39}%
}\hskip.1em\allowbreak{} 的循环空间，即存在向量\allowbreak{}\hskip.1em \RB{29.240}{7.948}{0.939}{%
\RG{0.000}{-0.002}{10.9091}{mi}{11}\RG{10.048}{-0.002}{10.9091}{sy}{50}\RG{20.351}{-0.002}{10.9091}{mi}{86}%
}\hskip.1em\allowbreak{} ，使得

\[
\RB{192.843}{10.073}{4.309}{%
\RG{30.756}{8.834}{10.9087}{ex}{0}\RG{185.249}{8.834}{10.9087}{ex}{1}\RG{101.640}{4.502}{7.9701}{mi}{110}\RG{107.032}{4.502}{7.9701}{sy}{0}\RG{113.851}{4.502}{7.9701}{rm}{49}\RG{146.232}{4.502}{7.9701}{mi}{110}\RG{0.000}{-0.001}{10.9091}{mi}{86}\RG{11.818}{-0.001}{10.9091}{rm}{61}\RG{23.333}{-0.001}{10.9091}{mi}{76}\RG{35.756}{-0.001}{10.9091}{mi}{11}\RG{42.759}{-0.001}{10.9091}{mi}{59}\RG{47.603}{-0.001}{10.9091}{mi}{39}\RG{54.759}{-0.001}{10.9091}{rm}{40}\RG{59.002}{-0.001}{10.9091}{mi}{11}\RG{66.021}{-0.001}{10.9091}{rm}{41}\RG{70.263}{-0.001}{10.9091}{mi}{59}\RG{75.095}{-0.001}{10.9091}{mi}{58}\RG{79.939}{-0.001}{10.9091}{mi}{58}\RG{84.782}{-0.001}{10.9091}{mi}{58}\RG{89.614}{-0.001}{10.9091}{mi}{59}\RG{94.458}{-0.001}{10.9091}{mi}{39}\RG{118.744}{-0.001}{10.9091}{rm}{40}\RG{122.987}{-0.001}{10.9091}{mi}{11}\RG{130.006}{-0.001}{10.9091}{rm}{41}\RG{134.248}{-0.001}{10.9091}{mi}{59}\RG{139.080}{-0.001}{10.9091}{mi}{39}\RG{152.170}{-0.001}{10.9091}{rm}{40}\RG{156.413}{-0.001}{10.9091}{mi}{11}\RG{163.432}{-0.001}{10.9091}{rm}{41}\RG{167.673}{-0.001}{10.9091}{mi}{59}\RG{172.506}{-0.001}{10.9091}{mi}{58}\RG{177.349}{-0.001}{10.9091}{mi}{58}\RG{182.182}{-0.001}{10.9091}{mi}{58}\RG{190.248}{-0.001}{10.9091}{mi}{58}%
}
\]
解：由秩–零度定理，dimker\allowbreak{}\hskip.1em \RB{26.735}{7.767}{2.876}{%
\RG{0.000}{0.002}{10.9091}{mi}{39}\RG{10.167}{0.002}{10.9091}{rm}{61}\RG{21.664}{0.002}{10.9091}{rm}{49}%
}\hskip.1em\allowbreak{} 。设\allowbreak{}\hskip.1em \RB{9.751}{5.324}{0.618}{%
\RG{0.000}{0.002}{10.9091}{mi}{109}%
}\hskip.1em\allowbreak{} 为\allowbreak{}\hskip.1em \RB{7.252}{5.324}{2.876}{%
\RG{0.000}{0.002}{10.9091}{mi}{39}%
}\hskip.1em\allowbreak{} 的幂零指数，即

\[
\RB{111.269}{10.077}{2.876}{%
\RG{7.136}{4.507}{7.9701}{mi}{109}\RG{69.223}{4.507}{7.9701}{mi}{109}\RG{76.965}{4.507}{7.9701}{sy}{0}\RG{83.782}{4.507}{7.9701}{rm}{49}\RG{0.000}{0.002}{10.9091}{mi}{39}\RG{18.452}{0.002}{10.9091}{rm}{61}\RG{29.950}{0.002}{10.9091}{rm}{48}\RG{35.422}{0.002}{10.9091}{mi}{59}\RG{62.072}{0.002}{10.9091}{mi}{39}\RG{91.706}{0.002}{10.9091}{sy}{54}\RG{91.706}{0.002}{10.9091}{rm}{61}\RG{103.204}{0.002}{10.9091}{rm}{48}\RG{108.675}{0.002}{10.9091}{mi}{58}%
}
\]
选择\allowbreak{}\hskip.1em \RB{29.240}{7.954}{0.933}{%
\RG{0.000}{0.003}{10.9091}{mi}{11}\RG{10.048}{0.003}{10.9091}{sy}{50}\RG{20.351}{0.003}{10.9091}{mi}{86}%
}\hskip.1em\allowbreak{} ，使\allowbreak{}\hskip.1em \RB{62.137}{9.532}{3.224}{%
\RG{0.000}{0.003}{10.9091}{mi}{39}\RG{7.136}{3.962}{7.9701}{mi}{109}\RG{14.876}{3.962}{7.9701}{sy}{0}\RG{21.694}{3.962}{7.9701}{rm}{49}\RG{26.588}{0.003}{10.9091}{rm}{40}\RG{30.831}{0.003}{10.9091}{mi}{11}\RG{37.849}{0.003}{10.9091}{rm}{41}\RG{45.121}{0.003}{10.9091}{sy}{54}\RG{45.121}{0.003}{10.9091}{rm}{61}\RG{56.619}{0.003}{10.9091}{rm}{48}%
}\hskip.1em\allowbreak{} 。

第一步：证明链长\allowbreak{}\hskip.1em \RB{9.751}{5.326}{0.616}{%
\RG{0.000}{0.004}{10.9091}{mi}{109}%
}\hskip.1em\allowbreak{} 不超过\allowbreak{}\hskip.1em \RB{6.729}{5.326}{0.616}{%
\RG{0.000}{0.004}{10.9091}{mi}{110}%
}\hskip.1em\allowbreak{} 。若有非平凡关系

\[
\RB{79.992}{19.205}{14.930}{%
\RG{0.000}{13.634}{7.9701}{mi}{109}\RG{7.741}{13.634}{7.9701}{sy}{0}\RG{14.560}{13.634}{7.9701}{rm}{49}\RG{1.576}{10.362}{10.9087}{ex}{88}\RG{37.179}{4.502}{7.9701}{mi}{106}\RG{20.725}{-0.002}{10.9091}{mi}{99}\RG{30.042}{-0.002}{10.9091}{mi}{39}\RG{41.775}{-0.002}{10.9091}{rm}{40}\RG{46.018}{-0.002}{10.9091}{mi}{11}\RG{53.037}{-0.002}{10.9091}{rm}{41}\RG{60.291}{-0.002}{10.9091}{rm}{61}\RG{71.800}{-0.002}{10.9091}{rm}{48}\RG{77.278}{-0.002}{10.9091}{mi}{59}\RG{25.447}{-1.638}{7.9701}{mi}{106}\RG{1.906}{-12.872}{7.9701}{mi}{106}\RG{5.957}{-12.872}{7.9701}{rm}{61}\RG{12.653}{-12.872}{7.9701}{rm}{48}%
}
\]
取满足\allowbreak{}\hskip.1em \RB{29.596}{8.314}{3.690}{%
\RG{0.000}{0.003}{10.9091}{mi}{99}\RG{4.721}{-1.633}{7.9701}{mi}{106}\RG{12.461}{0.003}{10.9091}{sy}{54}\RG{12.461}{0.003}{10.9091}{rm}{61}\RG{24.078}{0.003}{10.9091}{rm}{48}%
}\hskip.1em\allowbreak{} 的最小下标\allowbreak{}\hskip.1em \RB{5.256}{5.325}{0.617}{%
\RG{0.000}{0.003}{10.9091}{mi}{114}%
}\hskip.1em\allowbreak{} ，两边作用\allowbreak{}\hskip.1em \RB{37.271}{9.532}{2.875}{%
\RG{0.000}{0.003}{10.9091}{mi}{39}\RG{7.136}{3.962}{7.9701}{mi}{109}\RG{14.876}{3.962}{7.9701}{sy}{0}\RG{21.694}{3.962}{7.9701}{rm}{49}\RG{26.043}{3.962}{7.9701}{sy}{0}\RG{32.861}{3.962}{7.9701}{mi}{114}%
}\hskip.1em\allowbreak{} 。下标小于\allowbreak{}\hskip.1em \RB{5.256}{5.325}{0.617}{%
\RG{0.000}{0.003}{10.9091}{mi}{114}%
}\hskip.1em\allowbreak{} 的系数均为零，下标大于\allowbreak{}\hskip.1em \RB{5.256}{5.325}{0.617}{%
\RG{0.000}{0.003}{10.9091}{mi}{114}%
}\hskip.1em\allowbreak{} 的项因\allowbreak{}\hskip.1em \RB{35.702}{7.829}{2.875}{%
\RG{0.000}{0.003}{10.9091}{mi}{39}\RG{7.136}{3.962}{7.9701}{mi}{109}\RG{18.566}{0.003}{10.9091}{rm}{61}\RG{30.184}{0.003}{10.9091}{rm}{48}%
}\hskip.1em\allowbreak{} 而消失，于是

\[
\RB{74.358}{10.073}{3.228}{%
\RG{16.688}{4.502}{7.9701}{mi}{109}\RG{24.429}{4.502}{7.9701}{sy}{0}\RG{31.247}{4.502}{7.9701}{rm}{49}\RG{0.000}{-0.001}{10.9091}{mi}{99}\RG{9.551}{-0.001}{10.9091}{mi}{39}\RG{36.140}{-0.001}{10.9091}{rm}{40}\RG{40.383}{-0.001}{10.9091}{mi}{11}\RG{47.402}{-0.001}{10.9091}{rm}{41}\RG{54.656}{-0.001}{10.9091}{rm}{61}\RG{66.165}{-0.001}{10.9091}{rm}{48}\RG{71.643}{-0.001}{10.9091}{mi}{59}\RG{4.721}{-1.637}{7.9701}{mi}{114}%
}
\]
矛盾。因此\allowbreak{}\hskip.1em \RB{100.250}{9.534}{3.223}{%
\RG{0.000}{0.004}{10.9091}{mi}{11}\RG{7.003}{0.004}{10.9091}{mi}{59}\RG{11.847}{0.004}{10.9091}{mi}{39}\RG{19.003}{0.004}{10.9091}{rm}{40}\RG{23.245}{0.004}{10.9091}{mi}{11}\RG{30.264}{0.004}{10.9091}{rm}{41}\RG{34.506}{0.004}{10.9091}{mi}{59}\RG{39.338}{0.004}{10.9091}{mi}{58}\RG{44.181}{0.004}{10.9091}{mi}{58}\RG{49.025}{0.004}{10.9091}{mi}{58}\RG{53.857}{0.004}{10.9091}{mi}{59}\RG{58.701}{0.004}{10.9091}{mi}{39}\RG{65.884}{3.964}{7.9701}{mi}{109}\RG{73.624}{3.964}{7.9701}{sy}{0}\RG{80.442}{3.964}{7.9701}{rm}{49}\RG{85.336}{0.004}{10.9091}{rm}{40}\RG{89.579}{0.004}{10.9091}{mi}{11}\RG{96.597}{0.004}{10.9091}{rm}{41}%
}\hskip.1em\allowbreak{} 线性无关，故\allowbreak{}\hskip.1em \RB{30.852}{7.442}{1.990}{%
\RG{0.000}{0.004}{10.9091}{mi}{109}\RG{12.608}{0.004}{10.9091}{sy}{20}\RG{24.123}{0.004}{10.9091}{mi}{110}%
}\hskip.1em\allowbreak{} 。

第二步：证明链长\allowbreak{}\hskip.1em \RB{9.751}{5.317}{0.624}{%
\RG{0.000}{-0.004}{10.9091}{mi}{109}%
}\hskip.1em\allowbreak{} 至少为\allowbreak{}\hskip.1em \RB{6.729}{5.317}{0.624}{%
\RG{0.000}{-0.004}{10.9091}{mi}{110}%
}\hskip.1em\allowbreak{} 。令\allowbreak{}\hskip.1em \RB{23.126}{7.946}{2.217}{%
\RG{0.000}{-0.004}{10.9091}{mi}{86}\RG{6.363}{-1.640}{7.9701}{mi}{107}\RG{14.761}{-0.004}{10.9091}{rm}{61}%
}\hskip.1em\allowbreak{} ker\allowbreak{}\hskip.1em \RB{11.965}{9.753}{2.882}{%
\RG{0.000}{-0.004}{10.9091}{mi}{39}\RG{7.136}{3.954}{7.9701}{mi}{107}%
}\hskip.1em\allowbreak{} （\allowbreak{}\hskip.1em \RB{25.638}{8.066}{1.999}{%
\RG{0.000}{-0.004}{10.9091}{mi}{107}\RG{9.053}{-0.004}{10.9091}{sy}{21}\RG{20.568}{-0.004}{10.9091}{rm}{49}%
}\hskip.1em\allowbreak{} ）。当\allowbreak{}\hskip.1em \RB{25.966}{8.066}{1.999}{%
\RG{0.000}{-0.004}{10.9091}{mi}{107}\RG{9.053}{-0.004}{10.9091}{sy}{21}\RG{20.568}{-0.004}{10.9091}{rm}{50}%
}\hskip.1em\allowbreak{} 时，\allowbreak{}\hskip.1em \RB{7.252}{5.317}{2.882}{%
\RG{0.000}{-0.004}{10.9091}{mi}{39}%
}\hskip.1em\allowbreak{} 将\allowbreak{}\hskip.1em \RB{11.192}{7.946}{2.217}{%
\RG{0.000}{-0.004}{10.9091}{mi}{86}\RG{6.363}{-1.640}{7.9701}{mi}{107}%
}\hskip.1em\allowbreak{} 映入\allowbreak{}\hskip.1em \RB{22.115}{7.946}{2.217}{%
\RG{0.000}{-0.004}{10.9091}{mi}{86}\RG{6.364}{-1.640}{7.9701}{mi}{107}\RG{11.185}{-1.640}{7.9701}{sy}{0}\RG{18.003}{-1.640}{7.9701}{rm}{49}%
}\hskip.1em\allowbreak{} ；其在\allowbreak{}\hskip.1em \RB{11.194}{7.946}{2.217}{%
\RG{0.000}{-0.004}{10.9091}{mi}{86}\RG{6.364}{-1.640}{7.9701}{mi}{107}%
}\hskip.1em\allowbreak{} 上的限制的核恰为ker\allowbreak{}\hskip.1em \RB{7.252}{5.321}{2.879}{%
\RG{0.000}{-0.001}{10.9091}{mi}{39}%
}\hskip.1em\allowbreak{} ，维数为\allowbreak{}\hskip.1em \RB{5.071}{7.764}{0.501}{%
\RG{0.000}{-0.001}{10.9091}{rm}{49}%
}\hskip.1em\allowbreak{} 。因此

\[
\RB{191.245}{9.291}{3.228}{%
\RT{0.000}{-0.001}{10.9091}{d}\RT{6.065}{-0.001}{10.9091}{i}\RT{9.098}{-0.001}{10.9091}{m}\RG{20.003}{-0.001}{10.9091}{mi}{86}\RG{34.764}{-0.001}{10.9091}{rm}{61}\RG{46.262}{-0.001}{10.9091}{rm}{49}\RG{54.138}{-0.001}{10.9091}{rm}{43}\RT{65.065}{-0.001}{10.9091}{d}\RT{71.130}{-0.001}{10.9091}{i}\RT{74.163}{-0.001}{10.9091}{m}\RG{85.069}{-0.001}{10.9091}{mi}{39}\RG{92.204}{-0.001}{10.9091}{rm}{40}\RG{96.446}{-0.001}{10.9091}{mi}{86}\RG{108.177}{-0.001}{10.9091}{rm}{41}\RG{115.449}{-0.001}{10.9091}{sy}{20}\RG{126.964}{-0.001}{10.9091}{rm}{49}\RG{134.840}{-0.001}{10.9091}{rm}{43}\RT{145.751}{-0.001}{10.9091}{d}\RT{151.816}{-0.001}{10.9091}{i}\RT{154.849}{-0.001}{10.9091}{m}\RG{165.754}{-0.001}{10.9091}{mi}{86}\RG{188.651}{-0.001}{10.9091}{mi}{58}\RG{26.366}{-1.637}{7.9701}{mi}{107}\RG{102.810}{-1.637}{7.9701}{mi}{107}\RG{172.118}{-1.637}{7.9701}{mi}{107}\RG{176.939}{-1.637}{7.9701}{sy}{0}\RG{183.757}{-1.637}{7.9701}{rm}{49}%
}
\]
由dim\allowbreak{}\hskip.1em \RB{30.855}{7.951}{2.137}{%
\RG{0.000}{-0.000}{10.9091}{mi}{86}\RG{6.364}{-1.637}{7.9701}{rm}{49}\RG{14.287}{-0.000}{10.9091}{rm}{61}\RG{25.785}{-0.000}{10.9091}{rm}{49}%
}\hskip.1em\allowbreak{} 归纳得到dim\allowbreak{}\hskip.1em \RB{32.316}{8.071}{2.213}{%
\RG{0.000}{-0.000}{10.9091}{mi}{86}\RG{6.363}{-1.637}{7.9701}{mi}{107}\RG{14.760}{-0.000}{10.9091}{sy}{20}\RG{26.275}{-0.000}{10.9091}{mi}{107}%
}\hskip.1em\allowbreak{} 。又\allowbreak{}\hskip.1em \RB{35.469}{7.826}{2.878}{%
\RG{0.000}{-0.000}{10.9091}{mi}{39}\RG{7.137}{3.958}{7.9701}{mi}{109}\RG{18.453}{-0.000}{10.9091}{rm}{61}\RG{29.951}{-0.000}{10.9091}{rm}{48}%
}\hskip.1em\allowbreak{} ，所以\allowbreak{}\hskip.1em \RB{38.083}{7.951}{2.213}{%
\RG{0.000}{-0.000}{10.9091}{mi}{86}\RG{6.363}{-1.637}{7.9701}{mi}{109}\RG{17.679}{-0.000}{10.9091}{rm}{61}\RG{29.194}{-0.000}{10.9091}{mi}{86}%
}\hskip.1em\allowbreak{} ，从而

\[
\RB{82.462}{9.292}{2.616}{%
\RG{0.000}{-0.000}{10.9091}{mi}{110}\RG{9.579}{-0.000}{10.9091}{rm}{61}\RT{21.093}{-0.000}{10.9091}{d}\RT{27.159}{-0.000}{10.9091}{i}\RT{30.191}{-0.000}{10.9091}{m}\RG{41.096}{-0.000}{10.9091}{mi}{86}\RG{58.775}{-0.000}{10.9091}{sy}{20}\RG{70.290}{-0.000}{10.9091}{mi}{109}\RG{79.868}{-0.000}{10.9091}{mi}{58}\RG{47.459}{-1.636}{7.9701}{mi}{109}%
}
\]
结合两步，\allowbreak{}\hskip.1em \RB{30.852}{5.322}{0.620}{%
\RG{0.000}{-0.000}{10.9091}{mi}{109}\RG{12.608}{-0.000}{10.9091}{rm}{61}\RG{24.123}{-0.000}{10.9091}{mi}{110}%
}\hskip.1em\allowbreak{} 。第一步得到的\allowbreak{}\hskip.1em \RB{6.729}{5.322}{0.620}{%
\RG{0.000}{-0.000}{10.9091}{mi}{110}%
}\hskip.1em\allowbreak{} 个向量便是一组基，故

\[
\RB{145.551}{10.075}{4.306}{%
\RG{30.756}{8.837}{10.9087}{ex}{0}\RG{134.248}{8.837}{10.9087}{ex}{1}\RG{101.641}{4.505}{7.9701}{mi}{110}\RG{107.033}{4.505}{7.9701}{sy}{0}\RG{113.851}{4.505}{7.9701}{rm}{49}\RG{0.000}{0.001}{10.9091}{mi}{86}\RG{11.818}{0.001}{10.9091}{rm}{61}\RG{23.333}{0.001}{10.9091}{mi}{76}\RG{35.757}{0.001}{10.9091}{mi}{11}\RG{42.760}{0.001}{10.9091}{mi}{59}\RG{47.604}{0.001}{10.9091}{mi}{39}\RG{54.760}{0.001}{10.9091}{rm}{40}\RG{59.002}{0.001}{10.9091}{mi}{11}\RG{66.021}{0.001}{10.9091}{rm}{41}\RG{70.264}{0.001}{10.9091}{mi}{59}\RG{75.096}{0.001}{10.9091}{mi}{58}\RG{79.940}{0.001}{10.9091}{mi}{58}\RG{84.772}{0.001}{10.9091}{mi}{58}\RG{89.616}{0.001}{10.9091}{mi}{59}\RG{94.459}{0.001}{10.9091}{mi}{39}\RG{118.745}{0.001}{10.9091}{rm}{40}\RG{122.987}{0.001}{10.9091}{mi}{11}\RG{130.006}{0.001}{10.9091}{rm}{41}\RG{142.957}{0.001}{10.9091}{mi}{58}%
}
\]
由于\allowbreak{}\hskip.1em \RB{33.119}{7.822}{2.881}{%
\RG{0.000}{-0.003}{10.9091}{mi}{39}\RG{7.136}{3.955}{7.9701}{mi}{110}\RG{16.104}{-0.003}{10.9091}{rm}{61}\RG{27.601}{-0.003}{10.9091}{rm}{48}%
}\hskip.1em\allowbreak{} ，继续加入后面的迭代向量不改变所张成的空间，原题结论成立。\allowbreak{}\hskip.1em \RB{8.365}{7.991}{0.503}{%
\RG{0.000}{-0.003}{10.9091}{am}{3}%
}\hskip.1em\allowbreak{}

小结：最长的幂零向量链天然线性无关。核只有一维，使各次幂的核每次最多增加一维，因此最长链必须有\allowbreak{}\hskip.1em \RB{6.729}{5.317}{0.624}{%
\RG{0.000}{-0.004}{10.9091}{mi}{110}%
}\hskip.1em\allowbreak{}项，并由此得到循环基。

整体总结

A07与A13都把“一个向量的连续像”组织成一组基；A08通过极化把二次条件化为双线性条件；A09–A12则把矩阵等式转成像、核及直和关系。遇到秩的等式，优先检查相关空间的包含关系与交空间。


\end{document}
